 % ============================================================
 % perry wasm paper - LaTeX source (XeLaTeX, CJK)
 % 由 docs/paper/perry-wasm-paper.md 机械转换生成；正文以 Markdown 为准。
 % 编译：xelatex main.tex （需 xeCJK + Noto CJK 字体；本目录不随附 PDF）
 % ============================================================
\documentclass[11pt, a4paper]{article}

\usepackage[margin=1in]{geometry}
\usepackage{amsmath,amssymb}
\usepackage{graphicx}
\usepackage{booktabs}
\usepackage{array}
\usepackage{float}
\usepackage{enumitem}
\usepackage{verbatim}
\usepackage{url}
\usepackage[colorlinks=true,linkcolor=black,citecolor=black,urlcolor=blue]{hyperref}

\usepackage{xeCJK}
\setCJKmainfont{Noto Serif CJK SC}[BoldFont=Noto Sans CJK SC Bold]
\setCJKsansfont{Noto Sans CJK SC}
\setCJKmonofont{Noto Sans Mono CJK SC}
\setmonofont{DejaVu Sans Mono}[Scale=0.85]

\setlength{\parindent}{0pt}
\setlength{\parskip}{0.5em}
\linespread{1.25}
\hypersetup{pdftitle={把 TypeScript 编译成 WebAssembly 并在无 JS 引擎的宿主上运行},
            pdfauthor={perry wasm demo 项目组}}

\title{把 TypeScript 编译成 WebAssembly 并在无 JS 引擎的宿主上运行\\[0.4em]
\large 一次从运行时模块化、性能归因到编译器上游修复的系统探索}
\date{2026-09-21}
\author{perry wasm demo 项目组}

\begin{document}
\maketitle
\thispagestyle{empty}






\noindent\textbf{材料来源：} 本仓库 \texttt{docs/}、\texttt{tools/}、\texttt{src/}、\texttt{host/}、\texttt{runtime-wasm/} 的一手实测记录；上游仓库 \href{https://github.com/PerryTS/perry}{PerryTS/perry}、\href{https://github.com/fn-a/typerry}{fn-a/typerry}、\href{https://github.com/bytecodealliance/wasm-micro-runtime}{bytecodealliance/wasm-micro-runtime}。


\bigskip\hrule\bigskip

\section{中文摘要}

TypeScript 的交付长期等同于交付源码：打包与压缩只增加阅读成本，不改变语义可读性。编译成原生二进制能解决源码外流，却把"一次编写、到处运行"换成"一次编写、到处编译"，交叉编译矩阵随平台数量增长。WebAssembly 提供了第三条路。本文记录一次完整的技术探索：把 perry 编译器产出的 TypeScript→wasm 模块跑在\textbf{没有任何 JavaScript 引擎}的宿主上，宿主只保留 WASI 的 \texttt{fd\_write}。


探索分三段。第一段解决"能不能跑"：把 perry 的 Rust 运行时以 \texttt{\#![no\_std]} 编译成独立的 wasm 模块，与业务模块由 WAMR 多模块机制链接，业务模块的 211 个 \texttt{rt.*} 导入全部由该模块提供，其中 13 个真实实现、198 个编译期生成的报错桩；正负向用例 6/6 通过，输出与 perry 自带 JS 宿主层逐字节一致。第二段解决"慢在哪里"：用同一份 \texttt{src/bench.ts}（fib(29) 加 10⁶ 次循环）建立六路对照，把总倍数按\textbf{乘性因子}分解为「引擎因子 × codegen 因子」并做闭合校验（误差 <5\%），用手写"干净对照 wasm"隔离引擎，用 V8 的 TurboFan 稳态与 QuickJS 旁证交叉验证。结论是引擎无过：干净 wasm 在 WAMR AOT 下与手写 C 同速（0.97×），根因是 perry wasm codegen 的类型擦除，\texttt{+} 与条件判定被强制过桥，桥函数体占 AOT 路耗时的 95.8\%。第三段解决"怎么修、修到多少"：先做零上游依赖的 wat 后处理 pass（4 程序 × 7 变体共 28 次逐字节一致、0 误判、bench 快 7.3×），再直接修改上游 \texttt{perry-codegen-wasm} 的发射点，结果是 \textbf{122.112 ms → 3.891 ms（快 31.4×）}，跨过等价手工特化的上界 17.185 ms，逼近全去盒的 3.325 ms。


\textbf{关键词：} WebAssembly；TypeScript 编译器；运行时模块化；性能归因；类型特化；WAMR


\bigskip\hrule\bigskip

\section{English Abstract}

TypeScript delivery has long meant source delivery. Bundling and minification raise reading cost but leave semantics legible. Compiling to a native binary solves source exposure but trades "write once, run anywhere" for "write once, compile everywhere", and the cross-compilation matrix grows with every platform. WebAssembly offers a third path. This paper reports a complete exploration: running perry's TypeScript-to-wasm module on a host with \textbf{no JavaScript engine at all}, where the host retains only WASI \texttt{fd\_write}.


The exploration has three stages. The first establishes feasibility: perry's Rust runtime is compiled as a standalone \texttt{\#![no\_std]} wasm module, linked with the application module through WAMR's multi-module mechanism; all 211 \texttt{rt.*} imports of the application are supplied by that module, 13 with real implementations and 198 with compile-time-generated trap stubs. All six demo steps pass, with byte-identical output against perry's bundled JavaScript host layer. The second stage locates the cost: six execution routes run the same \texttt{src/bench.ts} (fib(29) plus a 10⁶-iteration loop), and the total slowdown is factored \textbf{multiplicatively} into an engine factor and a codegen factor with closure checks (error below 5\%). A hand-written "clean" wasm isolates the engine; V8's TurboFan steady state and QuickJS provide independent cross-validation. The engine is not at fault: clean wasm under WAMR AOT runs at native speed (0.97×), while the root cause is type erasure in perry's wasm codegen, which forces \texttt{+} and conditions across a dynamic-dispatch bridge whose function body accounts for 95.8\% of the AOT route's runtime. The third stage asks how far a fix can go: a zero-upstream-dependency wat post-processing pass (28 byte-identical comparisons across 4 programs × 7 variants, zero false positives, 7.3× on the benchmark) is followed by a direct patch to the upstream \texttt{perry-codegen-wasm} emission sites, yielding \textbf{122.112 ms → 3.891 ms (31.4×)}, past the equivalence hand-specialization ceiling of 17.185 ms and close to the fully unboxed 3.325 ms.


\textbf{Keywords:} WebAssembly; TypeScript compiler; runtime modularization; performance attribution; type specialization; WAMR


\bigskip\hrule\bigskip

\section{1. 引言}

\subsection{1.1 交付 TypeScript 意味着交付源码}

TypeScript 编译成 JavaScript 之后仍然是可读的文本。打包、压缩、是否附带 source map，都只影响阅读的难易，不影响语义的暴露。对业务逻辑即核心资产的场景来说，这是一个长期无解的问题：产物即源码。


编译成原生二进制是一条彻底的出路。代价同样明确：产物与操作系统、CPU 架构、ABI 绑定，每个组合都要单独出一份，交叉编译矩阵随平台数量线性膨胀。对第三方分发的库或插件，"到处编译"的成本往往高于它换来的收益。


WebAssembly 是第三条路。产物是平台无关的字节码，只分发一次；运行只需要一个 wasm 运行时，于是"到处运行"由运行时侧提供，而不是由编译器侧穷举。


\subsection{1.2 两个研究问题}

perry 的 wasm 后端让这条路看上去可行，但它的产物形态立刻带出一个疑问。perry 输出的不是"裸算法模块"，而是一个完整程序：导出 \texttt{\_start} 与 \texttt{memory}，程序入口在启动时把全部业务逻辑跑完 [S1:31]。同时它声明了 \textbf{211 个 \texttt{rt.*} 导入}，覆盖字符串、console、Math、JSON、Date、Map/Set、Buffer、crypto 等运行时能力；\textbf{无论程序是否用到，实例化时全部要求可解析} [S1:71]。perry 自带的宿主层是一套 JavaScript 实现，其代码注释说明了设计取向：\emph{Runtime operations (strings, console, objects) are imported from JavaScript.} [S1:164]。


这决定了本文要回答的两个问题：


\begin{itemize}
\item \textbf{RQ1（源码保护）}：产物里是否只剩字节码，不含可读的源码语义；
\item \textbf{RQ2（到处运行还剩多少）}：这 211 个导入所代表的运行时语义，是必须由每个平台各自重写一遍，还是可以跟着 wasm 一起分发。
\end{itemize}

RQ2 是关键。如果运行时只能由宿主逐个实现，"到处运行"实际退化为"到处写一套运行时"，源码保护的好处也被稀释：分发物从源码换成了字节码，运行时那一摊照旧。RQ1 则相对简单，且答案带着折扣，见 3.1 与 8.4。


\subsection{1.3 本文的贡献}

\begin{enumerate}
\item \textbf{一条已跑通的架构路线}：把 perry 的 Rust 运行时按 wasm 目标编译成独立模块，与业务模块由 WAMR 多模块机制链接，宿主只剩 WASI 的一个调用（3.2–3.4）。\texttt{rt.*} 的调用约定、业务代码、codegen 均未改动。
\item \textbf{一套 \texttt{rt.*} ABI 的逆向方法}：在没有规范、但有可运行参考实现时，把参考实现当 oracle 插桩，而不是读源码猜（3.3）。
\item \textbf{一套可复用的性能归因方法学}：把总倍数拆成乘性因子、用干净对照 wasm 隔离引擎、做闭合校验、再用独立引擎交叉验证（第 5 章）。这套方法的产出不是"wasm 比原生慢 N 倍"，而是"这 N 倍里哪一部分属于引擎、哪一部分属于 codegen"。
\item \textbf{对自建基线的自审}：\texttt{gcc -O2} 的深度自内联使 1,664,079 次逻辑调用只发生 91,759 次真实 call，审计用 gdb 断点与 callgrind 双证纠正了这一口径错误（4.4）。
\item \textbf{两种修复与其定量验证}：零上游依赖的 wat 后处理 pass（6.2）与上游 codegen 发射点特化（6.3），后者把 AOT 路从 122.112 ms 降到 3.891 ms。
\end{enumerate}

\subsection{1.4 材料与方法声明}

本文不引用外部文献。全部素材来自本仓库的一手实测记录与上游公开仓库，属于"作者提供自有素材"的情形：方法学参照在既有材料上展开，不经过系统文献检索阶段，也不编造参考文献。所有数字均给出出处（\texttt{[S*:行号]}）与测量口径（样本数、统计量、扣除项、机器），凡原文标注为推断的数字，本文保留同等标注。


\bigskip\hrule\bigskip

\section{2. 背景}

\subsection{2.1 perry 与 typerry}

\href{https://github.com/PerryTS/perry}{perry} 是用 Rust 写的 TypeScript/JavaScript 编译器，前端走 SWC 解析，后端走 LLVM，主要产物是原生可执行文件 [S1:29]。它的 wasm 后端被抽出来单独发布为 npm 包 \href{https://github.com/fn-a/typerry}{\texttt{@typerry/node}}，内部链路是 SWC → 自研 HIR → wasm codegen，通过 napi-rs 暴露给 JavaScript，零运行时依赖 [S1:29]。


两条后端对"运行时"的处理方式不同，这也是本文全部问题的源头：


\begin{itemize}
\item \textbf{native 后端}：\texttt{crates/perry-runtime} 是 Rust 写的运行时（GC、JSValue、内置对象、字符串），\texttt{crates/perry-runtime-static} 把它打成 \texttt{libperry\_runtime.a}，\texttt{perry compile} 直接把这个静态库链进可执行文件 [S1:163]。运行时不是宿主契约，它就是源码。
\item \textbf{wasm 后端}：运行时操作被声明为对宿主的导入 [S1:164]。211 个导入和"每个宿主都得实现一遍"的问题由此而来。
\end{itemize}

本探索涉及的版本：WAMR 2.4.3，构建需 \texttt{WAMR\_BUILD\_MULTI\_MODULE=1}、\texttt{WAMR\_BUILD\_LIBC\_WASI=1}、\texttt{WAMR\_BUILD\_TARGET=X86\_64} [S3:16-17]；perry 预编译 release v0.5.1520（\texttt{perry-linux-x86\_64.tar.gz}，解压前 624,641,943 B、解压后 2.2 GiB，sha256 \texttt{3423d9fe…83d}）[S2:128-129]；上游 wasm 后端源码取自 vendored perry，commit \texttt{87ecb02b} [S7:5]。


\subsection{2.2 WAMR 与 WASI}

\href{https://github.com/bytecodealliance/wasm-micro-runtime}{WAMR}（WebAssembly Micro Runtime）是 C 实现的嵌入式 wasm 运行时，同时提供解释器（含 FAST\_INTERP 字节码翻译模式）、AOT 编译产物加载器，以及多模块链接能力。本探索用到它的三件事：多模块链接（一个模块的导入由另一个已注册模块提供）、AOT（\texttt{wamrc} 把 wasm 编成机器码）、WASI 支持。


WASI 在本探索里的作用被压到最小：运行时模块自己不带 libc，标准输出与错误只用到 \texttt{wasi\_snapshot\_preview1.fd\_write} 一个调用 [S1:188]。宿主只需调用 \texttt{wasm\_runtime\_set\_wasi\_args}，让运行时模块的 \texttt{fd\_write} 落到 stdio。


\subsection{2.3 相关工作定位}

本文不引外部文献，因此这里只说明技术坐标系。占位的做法有三种：把运行时语义留给宿主（perry wasm 后端的默认形态）、把运行时静态链进同一模块（perry native 后端，以及 wasm 侧尚未实施的"路线二"）、以及把运行时编成独立 wasm 模块由多模块机制链接（本文路线三）。Component Model 提供了第四种更干净的表达（用 WIT 声明 \texttt{rt} 接口），但它与本文这套"f64 位模式 + 线性内存槽位"的零拷贝约定冲突，且 WAMR 侧支持较弱，本探索未实施 [S1:171]。


\bigskip\hrule\bigskip

\section{3. 系统构造}

\subsection{3.1 perry wasm 产物形态}

先看产物本身。perry 的 wasm 是一个完整程序：


\begin{verbatim}
Export:  _start / memory / __indirect_function_table / ...
Import[211]:
func[0] sig=1 <rt.string_new> <- rt.string_new
func[1] sig=2 <rt.console_log> <- rt.console_log
...
\end{verbatim}

211 个导入全挂在 \texttt{rt} 模块上 [S1:63-69]。也就是说，宿主想让它跑起来，就必须把这 211 个函数全部提供出来，哪怕只有三个真的会被调用 [S1:73]。


RQ1 的部分答案是正面的。检查 \texttt{app.wasm} 的段表，\textbf{没有 name 自定义段}（只有 type / import / func / table / memory / global / export / element / datacount / code / data 共十一段），函数名和局部变量名都不在产物里；\texttt{strings} 能捞出来的标识符全是程序自己的字符串字面量 [S1:243]。泄漏的是字面量与导入名。而导入名这一项泄漏得并不轻：198 个桩的名字（\texttt{array\_new}、\texttt{json\_parse}、\texttt{fetch\_url} 等）直接来自导入段，顺带把"这个程序会碰哪些运行时能力"标了出来 [S1:241]。因此源码保护成立，但门槛只是从"打开源码"抬到"反编译一遍再读"。


\subsection{3.2 运行时模块化：三条路线}

前文提到的三条路线，按改动量排序 [S1:168-172]：


\begin{center}\small\begin{tabular}{lrrr}
\toprule
路线 & 做法 & 状态 & 代价 \\
\midrule
一·C 桥接 & 宿主里用 C 手写 \texttt{rt.*}，编成 \texttt{libperry\_rt.so}，\texttt{iwasm --native-lib} 动态载入 & 已实施为探针，后废弃为历史背景 & 成本随程序用到的语言特性线性增长 \\
二·AOT 内联 & \texttt{wasm-ld} 把运行时 wasm 静态库与 codegen 输出链成单模块（同构于 native 路径）；或走 Component Model & 未实施 & 需 codegen 产出可重定位对象，或链接器把 import 解析成本地符号；Component Model 则与零拷贝约定冲突 \\
三·运行时 wasm 模块 & 运行时编成独立 wasm 模块，导出同名 \texttt{rt.*} 与 memory，业务模块 import 它，WAMR 多模块链接 & \textbf{当前实现，已实测} & 运行时的 OS 强耦合模块需裁剪（见 3.5） \\
\bottomrule
\end{tabular}\end{center}


路线一是探针而非终点，它把 ABI 测了出来，但也证明了"每个平台重写一遍运行时"不是答案：13 个手写实现只够支撑纯原始值的程序，换成用数组的程序立刻报错 [S1:23]。


路线三的构件如下 [S1:182-184, S3:38-53]：


\begin{itemize}
\item \texttt{runtime-wasm/}：Rust \texttt{\#![no\_std]} 运行时，\texttt{crate-type = ["cdylib"]}，release 用 \texttt{opt-level = "s"} + \texttt{lto = true} + \texttt{panic = "abort"} + \texttt{codegen-units = 1} + \texttt{strip}；\texttt{src/lib.rs} 620 行；自定义 \texttt{panic\_handler} 指向 \texttt{trap()}。
\item \texttt{.cargo/config.toml}：\texttt{--global-base=2097152}，把运行时自己的 data/bss/stack 放到 2 MiB 以上，避开业务模块的低地址区。
\item \texttt{tools/patch-app-memory.mjs}：把业务模块的 memory 段改成 \texttt{import rt.memory}（默认 \texttt{min=1} 页），因为同一块线性内存只能有一个定义者 [S3:135-144]。
\item \texttt{tools/gen-rt-symbols.mjs}：解析 \texttt{app.wasm} 导入段（手写 uleb 解析 types/imports），源文件里已定义 \texttt{rt\_<名字>(} 的只登记进符号表，其余生成"调用即报错"的桩 [S3:113-133]。
\end{itemize}

生成的符号表是编译期产物，不是运行时查表：


\begin{verbatim}
build/rt_symbols.rs: 211 个 rt 导入 (已实现 13, 桩 198)
\end{verbatim}

桩被调用时写 stderr 实名报错后 trap，绝不静默返回假数据 [S1:209]。这是一个刻意选择：比起悄悄返回 \texttt{undefined} 让程序带着坏数据往下跑，实名报错至少能让边界可见 [S1:211-217]。


13 个真实实现覆盖 string（new/len/eq/concat/to\_string）、number 与 bool 的 NaN-box 值、plus、console（log/warn/error）、以及动态分派（\texttt{mem\_call}/\texttt{mem\_call\_i32}）[S5:77-79, S3:93-100]。


\textbf{为什么这条路走得通}：211 个 \texttt{rt.*} 签名只用到 i32/i64/f32/f64，指针就是业务线性内存偏移，跨模块不存在"类型不匹配"这一层；运行时模块 import 业务那块内存后，\texttt{mem\_call}/\texttt{string\_new} 直接读写同一块内存，零拷贝约定原样成立 [S1:182]。


\subsection{3.3 \texttt{rt.*} ABI 的逆向：把参考实现当 oracle}

\texttt{rt.*} 的调用约定是 perry 内部约定，没有文档 [S1:77]。但 perry 自带一套宿主层，那 211 个函数的定义是现成的。


方法很简单：在宿主的 \texttt{mem\_call} 实现里插一行打印 [S1:81-85]。


\begin{verbatim}
const coreFn = __memDispatch[name];
// 上面插一行:
process.stderr.write(`MEMCALL ${name} args=${JSON.stringify(args)}\n`);
\end{verbatim}

插桩立刻给出事实：


\begin{verbatim}
MEMCALL js_add args=[6764,4181]
MEMCALL js_add args=["fib(0..19) sum = ",10945]
MEMCALL console_log args=["fib(0..19) sum = 10945"]
\end{verbatim}

一个 20 行程序真正用到的导入只有 \texttt{string\_new}、\texttt{mem\_call}、\texttt{mem\_call\_i32} 三个；所有动态调用（字符串相加、console 输出、\texttt{.length}）都收敛到 \texttt{mem\_call} 这一个入口，剩下 200 多个导入在实例化时被解析但一辈子不会被调用 [S1:89-95]。


方法论可以直接写成一句话：\textbf{当你没有规范、但有一个能跑的参考实现时，不要读源码猜，去插桩打印}。最初尝试从 wat 反推参数含义，盯着 \texttt{i64.const 9223090561878065386} 看了半小时没有结论；插一行 print 两分钟就清楚了 [S1:97]。


由此得到的 ABI 事实：


\begin{itemize}
\item \textbf{值编码是 NaN-boxing}，i64 里装 f64 位模式，高 16 位是标签 [S1:103-111, S3:60-66]：
\end{itemize}

\begin{center}\small\begin{tabular}{lr}
\toprule
值 & 位模式 \\
\midrule
\texttt{undefined} / \texttt{null} / \texttt{false} / \texttt{true} & \texttt{0x7FFC…0001} \textasciitilde{} \texttt{0x7FFC…0004} \\
对象/数组/闭包（handle） & 高 16 位 \texttt{0x7FFD}，低 32 位是 handle id \\
int32 快路径 & 高 16 位 \texttt{0x7FFE} \\
字符串 & 高 16 位 \texttt{0x7FFF}，低 32 位是字符串表下标 \\
其他 & 就是普通 double \\
\bottomrule
\end{tabular}\end{center}


\begin{itemize}
\item \textbf{字符串表是隐式契约}：wasm 启动时按固定顺序逐个调用 \texttt{rt.string\_new(offset, len)} 注册字面量，宿主必须按同样顺序 append，下标即 id；两边计数错一位，所有字符串就全乱了。桥接函数名也在这张表里，\texttt{mem\_call} 的第一个参数就是名字的下标 [S1:113]。运行时侧的表是 \texttt{MAX\_STRINGS=1024} 条、\texttt{ARENA\_SIZE=64 KiB}，条目为 \texttt{\{ptr, len, utf16\}}，其中 \texttt{utf16} 是 JS \texttt{length} 所需的 UTF-16 码元数 [S3:72-80]。
\item \textbf{动态调用协议是 \texttt{mem\_call(nameId, argc, base)}}：参数以 u64 槽位写在业务线性内存 \texttt{base} 处，返回值也写回 \texttt{base}（函数本身返回 0.0 当占位，宿主层源码注释说明了这点）；\texttt{mem\_call\_i32} 一样，只是结果直接作为 i32 返回、不写内存 [S1:115]。至于为什么不按 f64 直接传参而要绕一趟内存，源码没有解释；本文作者的判断是 f64 过 FFI 边界时 NaN 位模式有被规范化的风险，两边都按 u64 读写原始位模式最稳妥 [S1:115]。这条判断属于推断，未在文档中给出依据。
\end{itemize}

\subsection{3.4 多模块链接与 WASI-only 宿主}

\begin{verbatim}
graph LR
  A["src/app.ts"] -->|perry| B["build/app.wasm"]
  B -->|"import rt.memory + 211 个 rt.*"| C["build/rt.wasm<br/>Rust #![no_std] 运行时<br/>13 实现 + 198 桩"]
  C -->|"export rt.* + memory"| D["WAMR 多模块 runner<br/>host/perry_link.c"]
  D -->|"WASI fd_write"| E["stdout / stderr"]
\end{verbatim}

宿主 \texttt{host/perry\_link.c} 的顺序是 load rt → register → set\_wasi\_args → load app → instantiate，之后用 \texttt{wasm\_application\_execute\_main} 跑 \texttt{\_start} [S3:148-157]。\textbf{整个文件没有任何一行 \texttt{rt.*} 的实现}，这是与路线一最本质的区别 [S3:157]。


实测结果：\texttt{./demo.sh} 6/6 步 PASS，正向输出与 perry JS 宿主层逐字节一致（5 行），负向（用数组的 TS 程序）报 \texttt{bridge function 'array\_new' is not implemented}、退出码 1 [S1:186]。


\begin{verbatim}
fib(0..19) sum = 10945
Hello, WAMR!
msg.length = 12
string compare ok
template: Hello, WAMR! (sum=10945)
\end{verbatim}

\begin{verbatim}
Exception: bridge function 'array_new' is not implemented
execute _start: Exception: unreachable
\end{verbatim}

产物尺寸 [S3:254-260]：\texttt{app.wasm} 10650 B、\texttt{app\_link.wasm} 10658 B、\texttt{rt.wasm} 16798 B（采纳 nameId 缓存后为 16928 B，见 4.4）、宿主 runner \texttt{build/perry\_link} 531336 B。


\subsection{3.5 工程发现：WAMR AOT 不支持 import memory}

性能实验需要 AOT 引擎，于是撞上一个 WAMR 的文件格式限制：\textbf{AOT 产物不支持 import memory}。\texttt{core/iwasm/compilation/aot\_emit\_aot\_file.c} 硬编码 \texttt{import\_memory\_count = 0}（留了 TODO 注释），加载时 \texttt{core/iwasm/aot/aot\_validator.c} 直接拒绝（"import memory is not supported"）[S2:323-326]。


后果是：现有双模块结构（app import rt.memory）\textbf{无法整体 AOT}，\texttt{app.aot + rt.aot} 一起跑立刻 "out of bounds memory access" [S2:326-327]。混载方案（iwasm CLI 下 app.wasm 解释执行 + rt.aot）能跑，因为 AOT 子模块自身不 import memory，但业务代码仍被解释，对性能没有意义 [S2:329-330]。


可行的 AOT 形态是 \texttt{wasm-merge}（binaryen）把 app 与 rt 合并成单模块再交给 \texttt{wamrc}。合并产物有两个坑，靠后处理脚本 \texttt{tools/attribution/patch\_merged.mjs} 解决 [S2:330-334, S3:342-345]：


\begin{enumerate}
\item 合并产物导出了 rt 的 \texttt{\_\_data\_end}/\texttt{\_\_heap\_base}（1129776），与 app 的栈指针 global（65536）被 WAMR loader 组合成非法 aux stack（报 "auxiliary stack underflow"），必须删掉这两个导出；
\item 合并后 \texttt{\_initialize} 无人调用（原多模块形态由 WAMR 负责），必须织入 \texttt{\_start} wrapper 先调 \texttt{\_initialize}。
\end{enumerate}

还有一个副作用：合并单模块在多轮 \texttt{wasm\_application\_execute\_main} 下会在第 4 轮起报 "string table overflow"（rt 字符串表跨轮累积），解释器跑同一合并模块同样复现。这是合并形态的多轮限制，双模块原形态没有这个问题，因此 E 路计时改用每轮独立进程 [S2:62-67, S2:482-485]。


\subsection{3.6 实施中记录在案的七个坑}

\begin{center}\small\begin{tabular}{lrrr}
\toprule
\# & 现象 & 根因与修法 & 标注 \\
\midrule
1 & \texttt{rt.string\_new} 按 \texttt{u32,u32} 声明与业务模块 \texttt{(i32,i32)} 签名对不上 & 曾需 \texttt{wasm-abis=3}；仓库现无该痕迹，签名对齐靠 Rust 类型本身（\texttt{string\_new} 用 \texttt{u32}，其余桥用 \texttt{i64} 传 NaN-box 值） & [S3:172-179]，[INFERENCE] \\
2 & \texttt{perry\_rt\_unimplemented} 在产物里找不到，桩调用链接失败 & Rust ≥1.70 起 cdylib 只导出 \texttt{pub} 的 \texttt{\#[no\_mangle]} 符号；修法 \texttt{\#[no\_mangle] pub extern "C" fn} & [S3:181-186] \\
3 & 实例化报 \texttt{failed to link import memory (rt, memory)}、退出码 1 & 业务模块 import 的 memory \texttt{min} 超过 WAMR 对运行时模块记录的有效初始页数；实测 min=1 能过，min=2\textasciitilde{}100（含 17/18/64）全部失败，故默认 \texttt{min=1} & [S3:188-194] \\
4 & 报 \texttt{initializing thread failed!} & 与 \texttt{wasm\_runtime\_set\_wasi\_args} 调用时机有关；本 WAMR 2.4.3 构建未开 WASI 线程支持，\textbf{未复现} & [S3:196-203]，条件性 \\
5 & 跑完输出后尾部多余空行，或初始化执行两遍 & \texttt{execute\_func("\_start")} 与 \texttt{wasm\_application\_execute\_main} 的差异；统一用后者 & [S3:205-209] \\
6 & 编译报多余右花括号 & 调试清理残留，删除即可 & [S3:211-214] \\
7 & shell \texttt{sed} 多行插入把函数体拆散 & \texttt{\textbackslash{}n} 转义被吃掉；改用脚本做精确字符串替换，之后机械修改一律走脚本 & [S3:216-220] \\
\bottomrule
\end{tabular}\end{center}


第 3 条与第 4 条共同说明一件事：多模块链接的约束不止在 wasm 语义层，WAMR 自身的模块记录与初始化顺序同样构成契约。


\bigskip\hrule\bigskip

\section{4. 性能评估方法学}

架构跑通之后，下一个问题必然是"代价是多少"。这一章的组织顺序刻意与常见的基准章节相反：先讲对照怎么设计、测量纪律如何，再给结果，最后讲\textbf{我们如何审计自己的基线}。最后一步的产出是一条影响最大的修正。


\subsection{4.1 基准程序与六路对照}

基准程序 \texttt{src/bench.ts} 只有 20 行：递归 \texttt{fib(29)} 加一次 10⁶ 次求和循环 [S1:1-20]。


\begin{verbatim}
function fib(n: number): number {
  if (n < 2) return n;
  return fib(n - 1) + fib(n - 2);
}
const N_FIB = 29;
const N_LOOP = 1000000;
const f = fib(N_FIB);
console.log("fib(" + N_FIB + ") = " + f);
let sum = 0;
for (let i = 0; i < N_LOOP; i++) { sum += i; }
console.log("sum = " + sum);
\end{verbatim}

选择它有两个理由。一是\textbf{调用密集}：fib(29) 产生 1,664,079 次逻辑调用，正好把"跨语言边界的调用成本"放到显微镜下。二是\textbf{结果可校验}：两端输出 \texttt{fib(29) = 514229}、\texttt{sum = 499999500000}，可以在计时之前先做逐字节一致性检查。


六路对照共享同一份源码与同一组常量 [S2:496-505]：


\begin{center}\small\begin{tabular}{lrr}
\toprule
路 & 宿主 / 引擎 & 说明 \\
\midrule
A & WAMR FAST\_INTERP & perry wasm 模块 + rt.wasm 双模块，解释执行 \\
B & Node V8 & perry wasm 模块 + perry 自带 JS 宿主层 \\
C & 手写原生 & \texttt{gcc -O2}，i64 实现（\texttt{tools/bench\_native.c}） \\
D & perry 原生 & TS → LLVM → 可执行文件 \\
E & WAMR AOT & \texttt{wamrc} O3，wasm-merge 合并单模块（rt 代码也进机器码） \\
F & QuickJS & Bellard qjs 解释执行同一算法的 JS 版本 \\
\bottomrule
\end{tabular}\end{center}


对照的设计意图是有层次的：C 给出"这块硬件能有多快"的地板，D 回答"perry 自己的两条后端差多少"，E 回答"换成 AOT 引擎后还剩多少差距"，B 与 F 提供\textbf{独立引擎}的参照系。F 尤其重要，它刻意绕开 wasm 与桥，只回答一个问题：一个纯解释器跑同样的算法需要多久。


\subsection{4.2 测量纪律}

测量口径集中在 \texttt{docs/performance.md} 的"测量口径"节 [S2:40-70]。要点如下：


\begin{itemize}
\item \textbf{环境}：AMD Ryzen 7 5800H（8C16T，max 4.47 GHz），24 GiB RAM，Linux 6.17.13-2-pve，gcc 11.4.0，Node v24.21.0，rustc 1.95.0，WAMR 2.4.3 [S2:42-43]。
\item \textbf{解释器配置经实证}（\texttt{.deps/wamr-build/CMakeCache.txt}）：\texttt{WAMR\_BUILD\_INTERP=1}、\textbf{\texttt{WAMR\_BUILD\_FAST\_INTERP=1}}、\texttt{WAMR\_BUILD\_AOT=0}、\texttt{WAMR\_BUILD\_JIT=0}、\texttt{WAMR\_BUILD\_MULTI\_MODULE=1}、\texttt{CMAKE\_BUILD\_TYPE=Release}，即 FAST\_INTERP 字节码翻译模式、无 AOT、无 JIT [S2:44-47]。
\item \textbf{AOT 配置}：同上加 \texttt{WAMR\_BUILD\_AOT=1} + \texttt{WAMR\_BUILD\_WITH\_CUSTOM\_LLVM=1}，链接系统 LLVM 14 供 JIT 内联 stub；wasm→机器码由 \texttt{wamrc} 完成（wamrc 2.4.3 用系统 LLVM 14 构建，默认 O3/znver3）[S2:48-51]。
\item \textbf{时钟}：A、C 用进程内 \texttt{clock\_gettime(CLOCK\_MONOTONIC)}（毫秒）；B 与全部冷启动用 bash 内建 \texttt{time}（\texttt{TIMEFORMAT=\%3R}，本机无 \texttt{/usr/bin/time}）[S2:52-53]。
\item \textbf{样本}：稳态每目标 11 轮\textbf{交错执行}（A B C × 11），丢弃第 1 轮 warmup，取 10 个样本，报 P50 与最小/最大值；冷启动各 5 次取中位数 [S2:54-55]。
\item \textbf{正确性先于计时}：A/B/C 输出逐字节一致性校验是 \texttt{bench.sh} 的第 2 步，排在全部计时之前 [S2:450-451]。
\end{itemize}

每条路的计时方式因其产物形态不同而不同，这本身就需要写清楚：


\begin{itemize}
\item \textbf{D 路}：perry 产物没有多轮入口，单次执行 \textasciitilde{}6 ms 又低于 bash \texttt{time} 分辨率，因此用 \texttt{bench\_perry\_wrap.c} 对整个进程 fork/exec + waitpid 计时。数字包含进程启动 \textasciitilde{}0.6 ms（空载基线 fork+exec+\texttt{/bin/true} P50 0.568 ms），口径对 D \textbf{偏保守}[S2:56-59, S3:309-311]。
\item \textbf{B 路}：计时含 node 进程启动，已减去 \texttt{node -e ''} 空跑基线（优化前 23 ms、优化后 24 ms）[S2:60-61]。
\item \textbf{E 路}：合并单模块的多轮限制（见 3.5）迫使其每轮独立进程，每进程内 1 次预热 + 1 次计时，12 进程弃第 1 轮取 11 样本；E′ 无状态，与 A 同为进程内 11 轮 [S2:62-67]。
\item \textbf{A 路}：\texttt{bench\_time} 内 \texttt{wasm\_application\_execute\_main} 可重复执行，计时行来自进程内时钟，实例化开销单独输出 \texttt{INIT\_MS} [S2:68-70]。
\end{itemize}

\subsection{4.3 结果}

优化采纳后的主批（稳态，毫秒）[S2:136-146]：


\begin{center}\small\begin{tabular}{lrrrr}
\toprule
目标 & P50 & 最小 & 最大 & ÷C 原生 \\
\midrule
E. WAMR AOT（perry wasm，合并单模块） & 146.829 & 143.097 & 158.074 & 104.4× \\
E′. WAMR AOT（干净 wasm 对照） & 1.362 & 1.317 & 1.409 & 0.97× \\
A. WAMR 解释器 & 2470.678 & 2378.375 & 2527.101 & 1757.2× \\
B. Node V8（扣除启动 24 ms） & 1280.000 & 1246.000 & 1311.000 & 910.4× \\
C. 原生 gcc -O2 & 1.406 & 1.357 & 3.092 & 1× \\
D. perry 原生（进程级，含启动） & 6.286 & 5.946 & 7.306 & 4.5× \\
\bottomrule
\end{tabular}\end{center}


优化前（同一台机器、同一口径，rt 为按名扫描版）[S2:150-161]：A = 4009.746 ms（3866.041 / 4499.434），B = 1239.000 ms，C = 1.471 ms；倍数 A = 2725.9×、B = 842.3×。前后对比：A 路 4009.746 → 2470.678 ms（−38\%），codegen 因子 79× → \textasciitilde{}49×，总倍数 2726× → 1757×。


六路总表（含后增的 F 路）[S2:496-505]：


\begin{center}\small\begin{tabular}{lrrr}
\toprule
路 & P50 & ÷ C 原生 & ÷ D perry 原生 \\
\midrule
A. wasm × WAMR 解释器 & 2470.678 ms & 1757× & 393× \\
B. wasm × V8 & 1280.000 ms & 910× & 204× \\
C. 手写原生 gcc -O2 & 1.406 ms & 1× & — \\
D. perry 原生（进程级，纯执行 \textasciitilde{}1.6–2.6 ms） & 6.286 ms & 4.5×（纯执行 1.2–1.9×） & 1× \\
E. wasm × WAMR AOT（复测批 138.7 / 141.4） & 146.829 ms & 104× & 23×（vs D 纯执行 54–88×） \\
F. QuickJS 直跑 JS（进程级，qjs 空载 \textasciitilde{}0.95 ms） & 85.532 ms & 61× & 14×（进程级） \\
\bottomrule
\end{tabular}\end{center}


F 路的拆分与倍数 [S2:642-655]：P50 85.532 ms，其中 fib 部分 61.7 ms（37 ns/层）、loop 部分 24.8 ms（24.8 ns/迭代）；F/E = \textbf{0.58×}（F 更快）、F/E′ ≈ 62×、F/A = 0.035×（F 快 29×）。


冷启动（进程启动到输出完成，real 中位数 / 最小，毫秒）[S2:167-176]：A 4879 / 4836（\texttt{bench\_time} 固定先跑 1 次预热再计时，real 含两次执行；单次 ≈ 4879 − 2471 − 7 ≈ 2401 ms），A 的 \texttt{INIT\_MS} 6.883 / 8.042（双模块加载/注册/实例化，不含计算），B 1224 / 1202，C 3 / 3，D 6 / 6。优化前 A real 7836 / 7685（单次 ≈ 3819 ms），\texttt{INIT\_MS} 6.728 / 6.344。


\subsection{4.4 对自建基线的审计}

结果公布后出现一条质疑：原生基线 1.471 ms 在物理上可疑，因为按 1,664,079 次调用摊算，每层递归只有 0.6–0.9 ns，不可能。这条质疑指向的正是基准方法本身。


审计用三条独立证据核查 [S2:178-222]。


\textbf{证据一：多编译器交叉验证。} 同一 \texttt{bench\_native.c} 在 gcc -O1/-O2/-O3/-Ofast/-Os 与 clang -O2 下 P50 分别为 3.604 / 1.471 / 1.240 / 1.296 / 2.337 / 2.170 ms。五种独立编译管线聚在 1.2–3.6 ms 同一数量级；若有病态折叠，应出现离群值。


\textbf{证据二：指令计数闭合检验}（callgrind 实测）[S2:198-211]：


\begin{center}\small\begin{tabular}{lrrrr}
\toprule
编译 & 单次执行指令数 & 静态 call 点 & P50 & IPC（@3.0 GHz） \\
\midrule
gcc -O2（基线） & 16.71 M & 4（2 克隆体） & 1.471 ms & \textasciitilde{}3.8 \\
gcc -O1（对照） & 28.13 M & 2 & 3.604 ms & \textasciitilde{}2.6 \\
\bottomrule
\end{tabular}\end{center}


IPC 3.8 对简单整数短依赖链合理。若 gcc 跨 \texttt{printf} 合并了两次 fib 调用，每轮指令增量会减半，而实测每轮增量恒为 16.71 M；\texttt{fib(29)} 的 1,664,079 次逻辑调用由计数器实证（\texttt{fib\_count.c}，gcc -O2 下 count = 1,664,079，输出不变）。


\textbf{证据三：拆分计时}（\texttt{bench\_split.c}，gcc -O2）：fib 部分 P50 0.84 ms + 循环部分 0.30 ms ≈ 1.13 ms，与整体 1.47 ms 吻合（差额为两次 \texttt{clock\_gettime} 与 \texttt{printf}）。同型 f64 变体（\texttt{bench\_f64.c}，与 perry NaN-box 的 f64 位型同语义）P50 2.52 ms；C 基线用 i64 是公平下界，f64 版慢 \textasciitilde{}1.8×，量级不变。


\textbf{结论：计时数字成立，物理矛盾来自 \texttt{gcc -O2} 对 fib 的深度自内联。} 1,664,079 次"逻辑调用"只发生 \textbf{91,759 次真实 call}，由两个独立实测互相印证：gdb 断点在 warmup+1 个 RUN 上命中 183,519，除以 2 次顶层执行得 91,759；callgrind 调用图上 fib 与内联克隆体 fib'2 的入口合计同为 183,519。\texttt{fib\_O2.asm} 显示 161 条指令内含 4 个 call 点与 2 个克隆体，平均一次真实 call 覆盖约 18 层逻辑调用 [S2:182-188]。按真实 call 折算：fib 部分 0.84 ms / 91,759 ≈ \textbf{9.2 ns/真实 call}（≈27 cycles @3 GHz）。


质疑者的物理直觉对"真实 call"完全正确，错在把逻辑调用数当成了真实调用数。这条修正的意义超出这条基线本身：\textbf{原生基线执行的动态指令量远少于 wasm 路径在同语义下的指令量}，"1757×"里有一部分是代码形态差异，而不是全部都由解释器与桥的运行时代价构成。归因矩阵已把这一点计入引擎因子的分母一侧（干净 wasm × WAMR 35×），无需改数，但在呈现上必须按因子乘积来读 [S2:189-192]。


审计同时记录了两条来自开发过程的教训 [S2:446-460]：


\begin{enumerate}
\item \textbf{编辑器事故}：基准开发中 \texttt{src/bench.ts} 曾丢失 \texttt{const f = fib(N\_FIB)} 一行，导致 WAMR 路打印 \texttt{fib(29) = undefined}，排查中先排除了 rt.wasm 桩表与 codegen 路径，最后确认是源码编辑问题。教训是输出一致性校验必须先于计时。
\item \textbf{\texttt{gcc -O2} 会把整个基准常量折叠}：初版 \texttt{bench\_native.c} 直接用 \texttt{\#define} 常量，实测 0.002 ms（折叠后只剩 \texttt{printf}）；改经 \texttt{volatile} 指针读入 \texttt{N\_FIB}/\texttt{N\_LOOP} 后得到真实的 1.5 ms。\textbf{原生基线必须反折叠，否则倍数会虚高三个数量级。}
\end{enumerate}

\subsection{4.5 本轮审计的修正汇总}

性能数字不是一次成型。2026-09-21 的审计批针对三处质疑逐项复现，并新增 F 路，修正如下 [S2:663-678]：


\begin{center}\small\begin{tabular}{lrrr}
\toprule
项 & 旧值 & 新值 & 依据 \\
\midrule
B′ 干净 wasm × V8 & 4.9 ms & \textbf{3.40 ms}（\texttt{--no-liftoff}）；默认 4.67 仍对 & \texttt{run\_node\_steady.mjs} 阶梯预热 \\
B 路 V8 引擎因子 & \textasciitilde{}3.3× & \textbf{\textasciitilde{}2.5×} & 3.40 ÷ 1.471 \\
B 路 codegen+宿主层因子 & \textasciitilde{}253× & \textbf{\textasciitilde{}364×} & 1239 ÷ 3.40 \\
B 路乘积校验 & 3.3×253 ≈ 835 & \textbf{2.5×364 ≈ 910}（误差 <0.1\%） & 闭合 \\
D 纯执行 & 无（6.286 进程级） & \textbf{1.6–2.6 ms} & callgrind 19.44 M Ir 换算 \\
D/C（纯执行） & 4.5× & \textbf{1.2–1.9×} & 同左 \\
E/D & 23× & \textbf{54–88×}（按 D 纯执行） & 同左 \\
合并模块成本 & 未测 & \textbf{无额外成本} & 解释器跑合并 2340–2390 vs 双模块 2372–2480 ms \\
wamrc 优化级 & 声明默认 O3 & \textbf{确认默认 O3}，O0 对照 390 vs 141 & \texttt{wamrc --help} + 实测 \\
F 路 QuickJS & 无 & \textbf{85.532 ms}（进程级） & \texttt{bench\_f.sh} 11 轮 \\
E 路每桥成本 & \textasciitilde{}27.5 ns & \textbf{rt 函数体 25.4 ns + 装箱税 1.4 ns} & \texttt{nohost\_box} 隔离实验 \\
\bottomrule
\end{tabular}\end{center}


其中"合并模块无额外成本"这条决定 E 路的可信度：同一份合并模块喂解释器 iwasm，进程级 real 三测 2390 / 2381 / 2340 ms，双模块同口径 2480 / 2372 ms，\texttt{bench\_time} 双模块进程内 RUN 2501–2527 ms [S2:555-562]。合并版不慢于双模块，E 没有被合并本身高估。callgrind 在这里不可用：解释器构建含 \texttt{wrgsbase}（fsgsbase）指令，VEX 3.18 未实现，直接 SIGILL [S2:561-562, S3:372-373]。这是工具链对结论无关但在方法上必须记录的一处限制。


\bigskip\hrule\bigskip

\section{5. 归因方法学}

\subsection{5.1 问题：一个总倍数解释不了任何事}

A 路 2470.678 ms ÷ C 路 1.406 ms = 1757×。这个数字本身没有信息量：它把引擎代差、编译器产出的指令形态、桥接实现效率混成一团。若据此说"wasm 比原生慢 1757 倍"，读者会顺理成章地理解为"wasm 不行"，而真实情况可能是"某个编译器的某个后端没做特化"。归因的目标就是把这句话拆开。


\subsection{5.2 乘性分解与闭合校验}

方法分三层。


\textbf{第一层：把总倍数写成乘性因子的乘积。} 取一份与 \texttt{src/bench.ts} 完全同算法的\textbf{干净对照 wasm}（\texttt{tools/attribution/clean\_bench.wat}：纯 i64 指令、无 NaN-box、零 \texttt{rt.*} 导入，只有 \texttt{wasi fd\_write} 打印两行结果，零优化、手工书写）[S6:9]，把它跑在同一台 WAMR 上（A′，\texttt{clean\_time.c} 复用同一 \texttt{libiwasm.a}）和 node V8 上（B′，\texttt{run\_node.mjs}）。于是：


\[\text{总倍数} = \underbrace{\frac{\text{干净 wasm} \times \text{引擎}}{\text{原生}}}_{\text{引擎因子}} \times \underbrace{\frac{\text{perry wasm} \times \text{引擎}}{\text{干净 wasm} \times \text{引擎}}}_{\text{codegen 因子}}\]

\textbf{第二层：闭合校验。} 两个因子相乘必须能还原实测总倍数，否则说明还有第三个未被识别的成分。这是这套方法的硬约束，它把"讲故事"变成"可证伪"。


\textbf{第三层：独立引擎交叉验证。} 引擎因子与 codegen 因子都应该在换一个引擎后保持可解释：引擎因子应随引擎变化，codegen 因子应从"含解释器放大"退化为"纯机器码形态成本"。


\subsection{5.3 结果：三路矩阵}

\textbf{A 路（解释器）分解}，全部实测 [S2:267-277]：


\begin{center}\small\begin{tabular}{lrr}
\toprule
成分 & 倍数 & 证据 \\
\midrule
WAMR FAST\_INTERP vs 原生（干净代码下） & \textbf{\textasciitilde{}35×} & A′ 干净 wasm 50.8 ms（fib 43.0 + 循环 5.4）÷ 原生同形 1.47 ms（fib 0.967 + 循环 0.50） \\
perry codegen 差 + rt 桥实现 vs 干净 wasm（同在 WAMR） & \textbf{\textasciitilde{}49×} & A 2470.678 ms（优化后重测）÷ A′ 50.8 ms \\
\textbf{乘积校验} & 35×49 ≈ \textbf{1703} vs 实测 1757（误差 3.1\%） & 闭合 ✓ \\
\bottomrule
\end{tabular}\end{center}


优化前的同一矩阵 [S2:275-277]：codegen 因子 \textasciitilde{}79×（A 4010 ms ÷ A′ 50.8 ms），35×79 ≈ 2730 vs 实测 2726（误差 0.2\%）。79× → 49× 的差值就是"桥实现低效（按名线性扫描）"的贡献，已在 6.1 单独量化。


\textbf{E 路（AOT）分解}，全部实测 [S2:292-301]：


\begin{center}\small\begin{tabular}{lrr}
\toprule
成分 & 倍数 & 证据 \\
\midrule
WAMR AOT vs 原生（干净代码下） & \textbf{\textasciitilde{}0.97×}（与原生同速） & E′ 干净 wasm 1.362 ms ÷ 原生 1.406 ms \\
perry codegen 差 + rt 桥实现（同在 WAMR AOT，rt 代码也进机器码） & \textbf{\textasciitilde{}108×} & E 146.829 ms ÷ E′ 1.362 ms \\
\textbf{乘积校验} & 0.97×108 ≈ \textbf{105} vs 实测 104.4（误差 <1\%） & 闭合 ✓ \\
\bottomrule
\end{tabular}\end{center}


两个矩阵放在一起，事实就出来了 [S2:303-321]：


\begin{enumerate}
\item \textbf{引擎因子归零。} A′ 50.8 ms → E′ 1.362 ms（37×），干净 wasm 在 WAMR AOT 下与 \texttt{gcc -O2} 原生同速。解释器那 35× 是纯引擎开销，与 perry 无关。
\item \textbf{codegen 因子从 49× 变成 \textasciitilde{}108×，不是变差。} 解释器下 49× 的分母（A′ 50.8 ms）含解释器对\textbf{所有}代码的放大（干净代码也被拖慢 35×），桥调用 rt 侧同样被拖慢；AOT 下分母与 rt 侧都是机器码，剩下的差距是"perry NaN-box/桥调用形态的机器码 vs 干净 i64 机器码"的纯 codegen 成本。\textbf{49× 里解释器放大成分约占一半，另一半（机器码形态成本）在 AOT 下全部保留。} 这个因子变化是分母定义改变的结果，两个因子不可直接相除比较，属推断 [S2:308-314]。
\item \textbf{B 路被反超。} E 146.8 ms 比 B（V8 + perry JS 宿主层）1280 ms 快 8.7×。同一份 perry wasm，WAMR AOT + wasm rt 桥比 V8 + JS 宿主桥快一个数量级。A 路慢是解释器的锅，不是"wasm 路线不行"。
\item \textbf{对 perry 的定位。} 换到 AOT 后 perry wasm 路与 perry 原生路（D 6.286 ms）差 \textbf{\textasciitilde{}23×}（解释器下 393×），且这 23× 几乎全是 codegen 桥调用形态（引擎已归零）。
\end{enumerate}

\textbf{B 路分解} [S2:336-347]：node V8 wasm 引擎 vs 原生（干净代码下）\textbf{\textasciitilde{}2.5×}（B′ 3.40 ms，\texttt{node --no-liftoff} 强制 TurboFan 全优）；codegen + JS 宿主层因子 \textbf{\textasciitilde{}364×}；2.5×364 ≈ 910 与实测 910.4 闭合（误差 <0.1\%）。审计前这套数字是 3.3× 与 253×。


\subsection{5.4 第一性原理隔离实验：23× 到底是谁的}

乘性分解给出了因子，但"codegen 因子 108×"仍然是黑箱。质疑随之而来，且提得很有道理：perry 原生路与 perry→wasm→AOT 路都经过 LLVM 级优化，差 23× 不合常理。


于是做了一组隔离实验（\texttt{tools/attribution/nohost\_*.wat}，直接调用版 + \texttt{call\_indirect} 版）[S2:619-628, S3:376-381]：


\begin{center}\small\begin{tabular}{lrr}
\toprule
变体 & 构造 & P50 \\
\midrule
\texttt{nohost} 直接调用版 & 与 perry fib 完全相同的调用图（每层 2 次跨模块调用），被调方是平凡 wasm 函数 & \textbf{1.371 ms} \\
\texttt{nohost\_box} 版 & 同上，但用 \texttt{call\_indirect} 防内联，被调方做最小的 NaN-box i64↔f64 往返 & \textbf{5.913 ms}（税 4.604 ms） \\
fib 纯机器码 & \texttt{fib\_only.aot}，纯 i64 递归 fib(29) & 1.309 ms \\
loop 纯机器码 & \texttt{loop\_only.aot}（LLVM 把 10⁶ 循环常量折叠） & \textasciitilde{}0.002 ms \\
rt 侧 \texttt{mem\_call} 函数体 & 余项 & 135.5 ms \\
\textbf{合计} &  & \textbf{141.4 ms}，闭合 ✓（误差 <0.1\%） \\
\bottomrule
\end{tabular}\end{center}


三条读数：


\begin{enumerate}
\item \textbf{调用图形态本身不是瓶颈。} 与 perry 完全相同的调用图加上平凡被调方，AOT 编译器把被调方全部内联吸收，1.371 ms 与 fib 纯机器码 1.309 ms 基本相等。
\item \textbf{即使强制"不可内联的间接调用 + 装箱往返"，也只到 5.9 ms。} 税 4.604 ms 对应 533 万次桥 × 1.4 ns/桥。
\item \textbf{E 的 141.4 ms 减去 5.9 ms 得 135.5 ms，全部是 rt 侧 \texttt{mem\_call} 函数体的执行成本}，即"NaN-box 解码 + nameId 直查 + tag 分派 + f64 算术 + 编码"。折算每桥成本：rt 函数体 ≈ 135.5 ms / 5,328,158 桥 ≈ \textbf{25.4 ns/桥}；装箱往返税 ≈ 4.604 ms / 3,328,158 桥（仅 fib）≈ \textbf{1.4 ns/桥}。\textbf{rt 侧代码在合并后也进机器码，占 E 的 95.8\%。}
\end{enumerate}

对照解释器路：A 路 rt 桥约 463 ns/次（1.2 µs/层 ÷ 2）对 AOT 的 25.4 ns，差 18×，量级自洽（引擎代差加上桥函数体从解释执行变为直接运行）[S2:607-608]。


结论：\textbf{23× 不是"两个优化器的差距"，而是"喂给优化器的输入形态"的差距}：类型化 IR 与类型擦除装箱字节码之间的差距。唯一收敛路径是让 wasm codegen 做类型特化。


\subsection{5.5 独立引擎交叉验证}

归因结论必须能在别的引擎上复现，否则无法排除"WAMR 特有现象"。


\textbf{V8 侧}：\texttt{run\_node\_steady.mjs} 做阶梯预热（0/25/100/500/2000 次）观察 tier-up 收敛 [S2:564-580]。


\begin{center}\small\begin{tabular}{lrr}
\toprule
跑法 & 稳态 P50 & 说明 \\
\midrule
node 默认 & 4.671 ms & 与文档 4.9 ms 一致；fib 单函数 tier-up 预算未耗尽即结束 \\
\texttt{node --no-liftoff}（强制 TurboFan） & 3.400 ms & V8 真优形态 \\
\texttt{node --liftoff-only}（纯 baseline） & 9.482 ms & 纯 Liftoff 上界 \\
\bottomrule
\end{tabular}\end{center}


这一步证伪了"稳态实际 \textasciitilde{}1.3 ms"的假设（TurboFan 也不到 AOT 的 1.36 ms），同时把 V8 引擎因子从 3.3× 修正为 2.5×。原因是 wamrc O3 对全模块做 LLVM + znver3 编译，连 fib 自递归都能内联（同样是 91,759 次真实 call）；V8 TurboFan 对热路径保守 tier-up，不做递归内联，短递归 fib 上确实不如 AOT 全模块 O3。


\textbf{QuickJS 侧}（F 路）是最强旁证 [S2:657-661]：


> 一个纯解释器（QuickJS 解释 JS，fib 每层 37 ns）比 WAMR AOT 执行 perry 装箱字节码（每层 2 桥 × 26.8 ns + fib 本体 ≈ 55 ns）还快。


QuickJS 是"无桥的慢解释器"，E 是"有桥的机器码"，桥税 25.4 ns/次已经超过 QuickJS 解释一层 fib 除调用外的全部开销。这也解释了 B 路（V8 JS 宿主桥 0.39 µs/层）与 F 路的差距：JS 宿主桥比 wasm 桥再慢一个数量级。


三个引擎、三种实现路径指向同一结论，本文据此认为归因是稳健的：\textbf{引擎无过，根因在 perry wasm codegen 的类型擦除。}


\bigskip\hrule\bigskip

\section{6. 两种修复与其验证}

\subsection{6.1 问题定位：可内联的算术被强制过桥}

先把"过桥"的范围修正清楚，避免把问题说过头。用 \texttt{wasm2wat} 反汇编加上游 codegen 源码（\texttt{crates/perry-codegen-wasm/src/emit/}）双重证实 [S2:365-387]：


\begin{center}\small\begin{tabular}{lrr}
\toprule
操作 & 路径 & 证据 \\
\midrule
\texttt{+}（js\_add） & \textbf{过桥}（\texttt{mem\_call}） & \texttt{literals\_vars.rs}：\texttt{BinaryOp::Add => emit\_memcall("js\_add")}，注释 "handles string+number etc."；bench.wasm 中 fib 每层 1 次，nameId=8 \\
\texttt{-} \texttt{*} \texttt{/} & \textbf{内联} f64.sub/mul/div（含 reinterpret 对） & 同文件 \texttt{\_ =>} 分支 \\
\texttt{<} \texttt{<=} \texttt{>} \texttt{>=} & \textbf{内联} f64.lt/le/gt/ge & \texttt{Expr::Compare} 数值分支 \\
if/while/for 条件 & \textbf{过桥}（\texttt{mem\_call\_i32}，is\_truthy） & \texttt{stmt.rs} L66-69 等，bench.wasm 中 nameId=12 \\
\texttt{===} / \texttt{==} & 过桥（js\_strict\_eq） & \texttt{Compare} 分支 \\
字符串操作、console & 过桥（本来就必须） & \texttt{calls.rs} / \texttt{strings\_json.rs} \\
\bottomrule
\end{tabular}\end{center}


bench.ts 热路径的 2 次/层桥调用就是 \texttt{js\_add}（加法）加 \texttt{is\_truthy}（条件判定），不是全部算术。而要判断这是"值模型的必然"还是"可修复的缺陷"，有三条证据 [S2:389-409]：


\begin{enumerate}
\item \textbf{类型信息存在，wasm 后端一点没用。} \texttt{crates/perry-codegen-wasm/Cargo.toml} 只依赖 perry-hir / perry-codegen-js / perry-dispatch，\textbf{不依赖 perry-codegen}。类型化 ABI（\texttt{typed\_abi.rs}）、i32 快路径（\texttt{expr/i32\_fast\_path.rs}）、\texttt{Type::Int32} 消费全在原生 LLVM 后端。
\item \textbf{HIR 有完整类型基础设施。} \texttt{perry-hir/src/types.rs} 定义 \texttt{Int32}（注释为 "optimization for known integers"），\texttt{lower\_types.rs} 能把 number 表达式推成 \texttt{Type::Number}，\texttt{analysis/value\_types.rs} 是完整值类型推断。TS 是静态类型语言，\texttt{let sum = 0; sum += i} 的类型可静态获知。
\item \textbf{int32 快路径编码三处都有解码路径，wasm 后端从不发射。} \texttt{PERRY\_BOX\_INT32}（\texttt{0x7FFE}）在 ABI（\texttt{perry\_abi.h}）、runtime（\texttt{JSValue::int32}）、JS 宿主（\texttt{INT32\_TAG}）三处都有解码路径，但 wasm emit 全目录 grep \texttt{0x7FFE} 零命中。wasm 后端函数签名统一为 \texttt{vec![ValType::I64; n]}（\texttt{compile.rs} L711-713），无任何类型特化签名。
\end{enumerate}

裁决是：\textbf{codegen 未做类型特化与内联是缺陷}，同一编译器家族的原生后端已实现同等特化，wasm 后端是功能缺口；但它不是 wasm 后端独有的 bug，而是"wasm 后端落后于原生后端"的整体状态。真正属于"统一设计"的部分只有"所有用户值一律 NaN-box f64 位型"这一保守值模型 [S2:406-409]。


桥实现本身的低效也单独做了量化 [S2:411-428]。正式 rt 的 \texttt{invoke()} 原本对 10 项 \texttt{BRIDGES} 逐项 memcmp（\texttt{lib.rs} L582-586），而 nameId 本来就是稳定整数索引：


\begin{center}\small\begin{tabular}{lrr}
\toprule
rt 版本 & A 路 P50 (ms) & codegen 因子（÷ A′ 50.8 ms） \\
\midrule
正式 rt.wasm（按名扫描，优化前） & 3959–4010 & \textasciitilde{}79× \\
rt\_fast 实验（nameId 缓存直查） & \textbf{2292} & \textbf{\textasciitilde{}45×} \\
正式 rt.wasm（nameId 缓存直查，已采纳） & \textbf{2470.678} & \textbf{\textasciitilde{}49×} \\
\bottomrule
\end{tabular}\end{center}


实验期降幅 1718 ms（−43\%），按约 533 万次桥调用均摊 ≈ 322 ns/次，即每次桥调用从约 500 ns 降到约 180 ns。正式产物重测值 2470.678 ms 与实验值 2292 ms 偏差 +7.8\%（在机器噪声范围内）。改动是最小 diff：新增 \texttt{static mut NAME\_CACHE: [u8; 64]}，\texttt{invoke()} 加 nameId < 64 的缓存直查 fast path，按名扫描命中时回填；产物尺寸 16798 → 16928 B（+130 B）[S3:283-287]。


\subsection{6.2 零上游依赖的修复：通用桥内联后处理}

\textbf{实验 A：能不能靠现成工具解决。} 对合并产物跑 binaryen \texttt{wasm-opt} 123 [S2:712-728]：


\begin{center}\small\begin{tabular}{lrrr}
\toprule
变体 & 关键 flags & P50 (ms) & 相对 E \\
\midrule
E（未优化基线） & — & 122.112 & 1× \\
A1 & \texttt{-O3}（默认内联阈值） & 135.4 / 123.7† & ≈1× \\
A2 & \texttt{-O3 --always-inline-max-function-size=10000} & 99.6 / 96.0† & 0.78–0.81× \\
A3 & \texttt{-O4 --always-inline-max-function-size=10000} & 101.4† & 0.82× \\
A4 & \texttt{--inlining-optimizing --always-inline-max-function-size=10000 --precompute-propagate --dce} & 101.3† & 0.82× \\
\textbf{A5} & \texttt{--inlining-optimizing --always-inline-max-function-size=5000 --precompute-propagate --dce} & \textbf{93.497}（min 92.5 max 94.5） & \textbf{0.766×} \\
\bottomrule
\end{tabular}\end{center}


†A1/A3/A4 与 A2 同批（E 基线 131.232 ms，n=11）；A5 与 E/B1/B2 同批（E 基线 122.112 ms）。


\textbf{能内联，不能折叠分派。} \texttt{wasm-dis} 确认 A2 中 \texttt{call \$142/\$143} 全部消失（模块从 6714 行 wat 变成 277922 行），整条 \texttt{mem\_call} + \texttt{invoke} 被强制内联进 fib/loop，字面 nameId/argCount 的常量传播成功。但内联体里仍残留 11 路 \texttt{br\_table}，其索引来自 \texttt{NAME\_CACHE} 的\textbf{运行时内存 load}（\texttt{i32.load8\_u offset=1051877+nameId}）；binaryen 没有内存常量传播，内存内容运行时才确定（miss 路径会写缓存），所以 switch 无法消掉，完整 miss 路径的字符串查找代码也原样留在内联体里。\texttt{f64.const 8/12} 这种"常量分派"只有常量本身可折叠，分派表不可折叠 [S2:730-738]。


把 nameId 折叠成直接调用桥实现是可能的，但没有工具能达成：nameId 到 op 的映射活在 rt 的数据表加 \texttt{br\_table} 结构里，不是可识别的调用边；手工改 wat 可行，但收益被夹在 A5（93.5）与 B1（71.6）之间，而且只省 dispatch，\texttt{js\_add}/\texttt{is\_truthy} 的通用实现（类型打标、字符串分支、结果 unbox）仍在，属推断，不值得做 [S2:740-745]。


\textbf{结论}：纯后处理能把 E 从 122.1 ms 压到约 93.5 ms（−23\%），但内联下来的是"大 switch 搬运代码"，到不了特化量级；零上游改动的收益上限就是约 93 ms。代价是体积从 21 KB wasm 变成 661 KB、AOT 产物从 74 KB 变成 1.35 MB [S2:786]。


\textbf{实验 B：先量出天花板。} 不真改上游，对 perry 原样字节码做\textbf{等价手工特化}（\texttt{spec\_patch.py} 从 \texttt{bench\_merged\_patched.wat} 生成），替换的正是 codegen 类型特化会发射的指令 [S2:750-779]：


\begin{center}\small\begin{tabular}{lrrrr}
\toprule
变体 & 改动 & P50 (ms) & 相对 E & 相对 E′ \\
\midrule
E（perry 原样） & — & 122.112 & 1× & 100× \\
\textbf{B1} & 热路径 \texttt{+} 内联 \texttt{f64.add}，\texttt{is\_truthy} 桥保留 & 71.612 & 0.586× & 58.9× \\
\textbf{B2} & B1 + \texttt{is\_truthy} 内联为 \texttt{i64.ne} 假盒比较（NaN-box 与影子栈纪律保留） & \textbf{17.185}（min 16.5 max 21.0） & \textbf{0.141×（快 7.1×）} & 14.1× \\
V3 & 纯 f64：无盒、无影子栈、全内联 & 3.325 & 0.027×（快 36.7×） & 2.7× \\
E′ & clean\_bench（纯 i64，同批） & 1.216 & 0.010×（快 100×） & 1× \\
\bottomrule
\end{tabular}\end{center}


等价性论证决定了这几个数字能不能当作"codegen 特化后的真实预期"：此程序里 \texttt{js\_add} 两侧恒为 number（number + number 的 JS \texttt{+} 就是 f64.add，perry 的 number 表示即裸 f64 位型）；\texttt{is\_truthy} 的输入恒为 \texttt{f64.lt} 产出的盒布尔（\texttt{TAG\_TRUE} 0x7FF8000000000004 / \texttt{TAG\_FALSE} 0x7FF8000000000003），\texttt{is\_truthy(盒布尔) ≡ i64.ne v TAG\_FALSE}。替换保持影子栈增减逐指令不变；打印路径的字符串拼接 \texttt{js\_add} 与 \texttt{console\_log} 桥原样保留 [S2:763-767]。


B2 的乘性分解 [S2:769-779]：


\begin{itemize}
\item \textbf{E − B2 ≈ 105 ms} 是桥调用本体（每层 2 桥 × 约 2.08 M 桥）。类型特化把这部分全部消灭。
\item \textbf{B2 − E′ ≈ 16 ms} 是 perry 的影子栈\textbf{内存纪律}（每个值经 global sp 存/取内存，fib 每层约 20 条辅助指令，1.66 M 层 × 约 10 ns）。这不是 NaN-box 的税：B2 的 i64↔f64 reinterpret 对在机器码层面是空操作，LLVM O3 会消除（属推断），box 本身近零成本；16 ms 是"值走内存不走寄存器"的调用纪律成本。
\item \textbf{E′ ≈ 1.2 ms} 是纯机器码（LLVM 深度内联 fib）。
\end{itemize}

V3（3.3 ms）与 E′（1.2 ms）的差是 LLVM 对 f64 与 i64 两种 fib 的内联/优化差异（属推断），两版都是"无盒无纪律"，不代表 perry 可控项。


\textbf{实验 C：把实验 B 的知识自动化。} B2 的成功依赖"人读过 \texttt{src/bench.ts} 才知道那里是 number"，这份类型知识在 perry 产物里已被 codegen 擦除，因此 B2 只是\textbf{上界估计器}，不是可用修复。于是实现了一个通用后处理 pass（\texttt{tools/attribution/bridge\_inline\_pass.mjs}，约 800 行 JS，wat→wat），让它自己从模块里恢复类型信息 [S2:813-829]。


pass 的值域有三格：\texttt{NUM}（原始 f64 位型 = JS number）／\texttt{BOOLBOX}（\texttt{TAG\_TRUE}/\texttt{TAG\_FALSE} 二值盒布尔）／\texttt{OTHER}。\textbf{NUM 的语义依据}：perry 的 number 表示就是裸 f64 位型（rt \texttt{encode(V::Num(n)) = n.to\_bits()}，其余值一律 NaN-box），所以"f64 域生产者 ⇒ number"；而 perry 自己的 codegen 对 \texttt{-}/\texttt{*}/\texttt{/} 就是\textbf{无条件}内联 \texttt{F64Sub/F64Mul/F64Div}，即 perry 本身已把 f64 域操作数当 number，pass 不引入 perry 未有的假设 [S2:866-872]。


抽象解释在每个函数内按语句序进行，控制流合并取保守并 [S2:894-904]：


\begin{verbatim}
js_add（nameId 8, argc 2）：
  (drop (call $mem_call (f64.const 8) (f64.const 2) BASE))
→ (i64.store BASE (i64.reinterpret_f64 (f64.add
     (f64.reinterpret_i64 (i64.load BASE))
     (f64.reinterpret_i64 (i64.load (BASE+8))))))

is_truthy（nameId 12, argc 1）：
  (call $mem_call_i32 (f64.const 12) (f64.const 1) BASE)
→ (i64.ne (i64.load BASE) (i64.const TAG_FALSE))
\end{verbatim}

两处改写都与 rt 侧实现逐位等价（\texttt{js\_add(Num,Num) = Num(a+b)}、\texttt{truthy(Bool(b)) = b}）。\textbf{number 条件保守回退}：JS truthiness 里 \texttt{0}/\texttt{-0}/\texttt{NaN} 均 falsy，非二值，不内联，仍走桥。其余桥（\texttt{console\_log}=4、\texttt{string\_concat}、\texttt{string\_len}=10、\texttt{js\_strict\_eq}=13 等）一律不动；nameId 语义来自 rt 固定桥表，与数据段字符串序一致（\texttt{id = 序 + 1}）[S2:900-904]。


泛化验证用 4 个程序 × 7 个变体 [S2:906-936]：


\begin{center}\small\begin{tabular}{lrr}
\toprule
程序 & 形态 & 用途 \\
\midrule
\texttt{bench.ts} & 纯 number 热循环 & 基准 \\
\texttt{probe\_str.ts} & 字符串密集：\texttt{s = s + "ab"}（×32）、\texttt{s.length}、\texttt{s === s}（×20 万）、\texttt{hits + 1}、\texttt{n > 8} & 验证\textbf{不误伤}本该走桥的字符串运算 \\
\texttt{probe\_mixed.ts} & 混合类型：\texttt{total + i}（×20 万）、\texttt{i === 199999}、\texttt{label + "!"} & 验证类型判定边界 \\
\texttt{probe\_nested.ts} & 跨函数：\texttt{dbl}/\texttt{acc\_upto}/主循环，返回值就是 js\_add 结果 & 验证\textbf{跨过程}类型推断 \\
\bottomrule
\end{tabular}\end{center}


覆盖率（静态桥调用点改写比例，pass 自报）：


\begin{center}\small\begin{tabular}{lrrr}
\toprule
程序 & 桥调用点 改写前→后 & 总覆盖率（保守 / \texttt{--closed-world}） & 明细 \\
\midrule
bench & 10 → 6 & 4/10 = \textbf{40\%} / 40\% & \texttt{js\_add} 2/6（另 4 处操作数是字符串，正确拒绝）、\texttt{is\_truthy} 2/2、\texttt{console\_log} 0/2 \\
probe\_str & 11 → 6 & 5/11 = \textbf{45\%} / 45\% & \texttt{is\_truthy} 4/4、\texttt{js\_add} 1/2（\textbf{拒绝的是 \texttt{s + "ab"}}）、\texttt{string\_len}/\texttt{string\_eq}/\texttt{console\_log} 0/5 \\
probe\_mixed & 7 → 4 & 3/7 = \textbf{43\%} / 43\% & \texttt{is\_truthy} 2/2、\texttt{js\_add} 1/2（\textbf{拒绝的是 \texttt{label + "!"}}）、\texttt{js\_strict\_eq}/\texttt{console\_log} 0/3 \\
probe\_nested & 7 → 4 & 3/7 = \textbf{43\%} / \textbf{6/7 = 86\%} & \texttt{is\_truthy} 2/2；\texttt{js\_add} 1/4（保守）→ \textbf{4/4}（closed-world） \\
\bottomrule
\end{tabular}\end{center}


性能与正确性 [S2:925-947]：


\begin{center}\small\begin{tabular}{lrrrrrrr}
\toprule
程序 & base & pass & pass(cw) & pass+wasmopt & pass+wasmopt(强) & pass+segue & pass+wasmopt+segue \\
\midrule
bench & 123.098 & \textbf{16.958} & 16.822 & \textbf{16.034} & 16.203 & 17.315 & 17.661 \\
probe\_str & 30.297 & \textbf{14.942} & 14.340 & \textbf{10.213} & 14.248 & 13.429 & 11.908 \\
probe\_mixed & 21.261 & \textbf{5.918} & 5.861 & \textbf{4.421} & 5.629 & 5.058 & 4.690 \\
probe\_nested & 0.521 & \textbf{0.280} & \textbf{0.044} & 0.236 & 0.282 & 0.233 & 0.226 \\
\bottomrule
\end{tabular}\end{center}


\textbf{正确性 28/28 逐字节一致}（每个程序 7 个变体，参照物是 perry 自带 JS 宿主层 \texttt{wasmBoot} 的 \texttt{run.mjs}，比对方式与 \texttt{demo.sh} 同）；误判清单为\textbf{空}。三个非显然读数：字符串密集程序也快 2.0×（30.3 → 14.9，靠 \texttt{is\_truthy} 4/4 全内联加 \texttt{hits + 1} 的 js\_add 内联，而字符串 \texttt{+} 与 \texttt{===} 一处没动）；混合程序快 3.6×（21.3 → 5.9）；probe\_nested 的保守与 closed-world 差 6.4×（0.280 vs 0.044 ms），差距\textbf{全部}来自"导出函数参数是否可当 number"。


最后一点揭示了这套方法的\textbf{天花板}，而且它是方法固有的、不是实现缺陷 [S2:943-948]：perry 把每个用户函数都导出（\texttt{\_\_wasm\_func\_N}），保守模式无法排除"宿主用字符串调它"，于是 \texttt{dbl(n) \{ return n + n \}} 的参数不可证；合并后的 AOT 模块实际是封闭世界（只有 \texttt{\_start} 一个入口），\texttt{--closed-world} 显式声明这一点后跨函数推断全部打通（js\_add 4/4）。\textbf{类型知识在模块里不可恢复时，pass 只能保守拒绝。} 防护规则保证了误判不会静默发生：只在两侧都可证 number（或输入可证是二值盒布尔）时改写；\texttt{TAG\_TRUE}/\texttt{TAG\_FALSE} 从被测模块自身推导而不是硬编码（实测本仓 rt 产物与 rt 源码常量不一致）；影子栈增减逐指令不变；每个变体都必须通过逐字节一致性验收，失配即 fail-fast [S2:950-959]。


另一条被排除的路是"调开关" [S2:961-969]：perry CLI / \texttt{@typerry/node} / 环境变量（\texttt{--target wasm|web}、\texttt{--minify}、\texttt{--fast-math}、\texttt{--march=*}、\texttt{--no-auto-optimize}、\texttt{PERRY\_TARGET\_CPU}、\texttt{PERRY\_PRECOMPILE}）产出的 wasm \textbf{字节完全相同}（md5 \texttt{af3e4dd7…}，9827 B）。wamrc 侧唯一有效的是 \texttt{--enable-segue}（配 \texttt{--target=x86\_64 --disable-llvm-jump-tables}）：122.14 → 99.84 ms（−18.3\%），仍 71× 于原生，且被本 pass 覆盖（桥没了，segue 也就没有收益）。其余开关无效或更差（\texttt{--opt-level=0} 灾难性 3.2×、\texttt{--enable-shared-heap} +29\%、\texttt{--enable-llvm-pgo} 因缺 \texttt{WAMR\_BUILD\_STATIC\_PGO=1} 无法闭环未验证）。


\subsection{6.3 上游 patch：codegen 发射点特化}

后处理 pass 解决了"不能改上游时怎么办"，但它带一个永久性缺陷：依赖 perry 产物的指令形态，perry 一升级 codegen 就可能失配（见 7.2）。真正的修复在源头。


patch 的对象是 vendored perry（commit \texttt{87ecb02b}，typerry \texttt{d13b5769} 的 submodule），目标是 \texttt{crates/perry-codegen-wasm}（无 LLVM 依赖），diff 规模 6 文件 / +487 −20 [S7:5-8]。


\textbf{新增保守类型事实}（新文件 \texttt{src/emit/type\_facts.rs}，419 行）：收集声明类型加轻量数据流，提供 \texttt{expr\_is\_number} / \texttt{expr\_is\_boolean} [S7:22, S7:68-85]。


\begin{itemize}
\item \textbf{number 判据}：\texttt{Number}/\texttt{Integer} 字面量；声明为 \texttt{Number}/\texttt{Int32} 的局部；\texttt{Update}（\texttt{++}/\texttt{--}）、\texttt{Unary Neg/Pos}、\texttt{Binary Sub/Mul/Div}（这些 perry 无条件 f64 内联，产物恒 f64 位型，perry 自身已当 number 处理，不引入新假设）；返回类型声明为 number 的函数调用；两侧都可证的 \texttt{+}（递归）。
\item \textbf{boolean 判据}：\texttt{Bool} 字面量；\texttt{Compare}（发射恒为 \texttt{If(Result I64)} 选 TAG\_TRUE/TAG\_FALSE）；\texttt{!x}；声明为 \texttt{Boolean} 的局部；返回布尔可证的调用。
\item \textbf{数据流补充}（这是 \texttt{let sum = 0} 能被证明的关键）：无注解 \texttt{let x = <init>} 且 init 可证同型则 x 进入候选，随后做\textbf{赋值敏感不动点}：x 的每个赋值点 RHS 都必须可证同型才保留（\texttt{Update} 恒 number，不破坏 number 候选，但会破坏 boolean 候选故拒绝）。从乐观初值单调递减，收敛后剩余候选在任意执行路径上取值都可证同型。\textbf{被闭包捕获或被函数 \texttt{captures} 捕获的 id 一律拒绝}。扫描用 perry-hir 的 \texttt{walker::walk\_expr\_children}（穷尽匹配，编译期强制覆盖所有 Expr 变体，不漏赋值点）。
\end{itemize}

\textbf{两个发射点特化} [S7:26-27, S7:31-66]：


\begin{verbatim}
// 1. 加法（emit/expr/literals_vars.rs, BinaryOp::Add）
if self.expr_is_number(left) && self.expr_is_number(right) {
    self.emit_expr(func, left);   F64ReinterpretI64;
    self.emit_expr(func, right);  F64ReinterpretI64;
    F64Add; I64ReinterpretF64;
} else { /* 原 emit_frame_begin(2) + store_arg×2 + emit_memcall("js_add", 2) */ }

// 2. 条件（emit/stmt.rs 的 if / while / do-while / for 4 处）
if self.expr_is_boolean(condition) {
    self.emit_expr(func, condition);          // 栈上盒布尔 i64
    I64Const(TAG_FALSE); I64Ne;               // → i32
} else { /* 原 emit_frame_begin(1) + store_arg + emit_memcall_i32("is_truthy", 1) */ }
\end{verbatim}

\textbf{等价性论证}（决定这些数字是不是"真实产物"而非"手工变体"）：perry 的 number 表示即裸 f64 位型（rt \texttt{encode(V::Num(n)) = n.to\_bits()}）。若运行时两侧确为 number，\texttt{js\_add(Num(a), Num(b))} 返回 \texttt{Num(a+b)}，与内联的 \texttt{i64.reinterpret\_f64(f64.add(f64.reinterpret\_i64(a), f64.reinterpret\_i64(b)))} \textbf{逐位相同}（含 NaN/±0/Inf 传播，f64 加法语义一致）。影子栈方面，原路径 \texttt{emit\_frame\_begin(2)} 推进 sp+16、\texttt{emit\_memcall} 内收 sp−16，净 0；特化路径完全不碰 sp，净效果一致。字符串分支：只有两侧\textbf{都可证} number 才内联，\texttt{string + anything} 恒走原桥，JS \texttt{+} 的字符串拼接语义保留。布尔条件方面，\texttt{expr\_is\_boolean} 只接受"恒产 TAG\_TRUE/TAG\_FALSE 二值盒布尔"的表达式，对这些值 \texttt{is\_truthy(盒布尔) ≡ (v != TAG\_FALSE)} 逐位等价。\textbf{number 条件保守回退}：0/−0/NaN 均 falsy，裸 i64 比较会误判，故一律回退原桥 [S7:41-64]。


辅助修复一处：\texttt{compile.rs} 的 globals init 循环（\texttt{:1319}）与 class 注册循环（\texttt{:1339}）补 \texttt{self.current\_mod\_idx = mod\_idx;}（原本缺失），per-module 类型事实会取错索引 [S7:25]。


\textbf{性能}（同口径 \texttt{build/aot\_time}，12 轮弃第 1 轮取 11 样本中位数）[S7:110-121, S2:1027-1040]：


\begin{center}\small\begin{tabular}{lrrr}
\toprule
变体 & P50 (ms) & 相对 E & 说明 \\
\midrule
E（perry 原样） & 122.112 & 1× & — \\
\textbf{patch 后（codegen 特化）} & \textbf{3.891}（min 3.673 / max 4.009，n=11） & \textbf{0.0319×（快 31.4×）} & 达到并超过 17 ms 目标 4.4× \\
B2（等价手工特化，保留影子栈纪律） & 17.185 & 0.141× & 实验 B 天花板 \\
V3（纯 f64，无盒无影子栈） & 3.325 & 0.027× & 全去盒上限 \\
E′（clean i64） & 1.216 & 0.010× & 硬件下限 \\
\bottomrule
\end{tabular}\end{center}


\textbf{正确性} [S2:1042-1048, S7:123-125]：\texttt{fib(29) = 514229}、\texttt{sum = 499999500000}；\texttt{./demo.sh} \textbf{6/6 PASS}（含负向 \texttt{array\_new} 报错）；3 个泛化探针 \texttt{probe\_\{str,mixed,nested\}.ts} 的 patch 后产物经完整 E 路（rt 桩 + 链接 + WAMR）输出与 perry JS 宿主层参照逐字节一致（\texttt{64/200000/long}、\texttt{n!/19999900000}、\texttt{323400}）。


\textbf{反汇编证据}（\texttt{--bare} 产物 \texttt{build/bench.wasm}，\texttt{wasm-dis}）[S7:127-142, S2:1050-1065]：


\begin{center}\small\begin{tabular}{lrr}
\toprule
指标 & patch 前 & patch 后 \\
\midrule
文件大小 & 9780 B & 9561 B \\
\texttt{call \$mem\_call}（js\_add 等） & 8 & 6 \\
\texttt{call \$mem\_call\_i32}（is\_truthy） & 2 & 0 \\
\texttt{f64.add} & 0 & 3 \\
\texttt{i64.ne} & 0 & 2 \\
\bottomrule
\end{tabular}\end{center}


fib 热路径上，\texttt{if (n < 2)} 由 \texttt{mem\_call\_i32(is\_truthy)} 变成 \texttt{i64.ne TAG\_FALSE}；\texttt{fib(n-1) + fib(n-2)} 与循环 \texttt{sum += i} 由 \texttt{mem\_call(js\_add)} 变成 \texttt{f64.add}。剩余 6 处 \texttt{mem\_call} 全是字符串拼接（\texttt{"fib(" + … + ") = " + f}、\texttt{"sum = " + sum}），正符合"可证 number 才内联"的设计边界。


\subsection{6.4 为什么上游 patch 比手工特化还快 4.4×}

这一处需要解释，否则容易误读。B2 是"手工等价特化"，patch 后产物在语义上做的是同一件事，为什么 3.891 ms 会远快于 B2 的 17.185 ms（快 4.4×）？


答案是\textbf{替换的粒度不同} [S2:1037-1040, S7:119-121]：B2 的手工替换只动 \texttt{mem\_call} 调用本身，其外围的\textbf{帧建立}与\textbf{影子栈内存槽往返}指令原样保留；而 codegen 发射点特化让整条帧建立与内存槽往返\textbf{都不再发射}，产物更紧，AOT 后端因此能更好优化。patch 后产物因此跨过 B2 天花板 17.185 ms，逼近 V3 的 3.325 ms。这条因果解释在原文中标注为推断（\texttt{[INFERENCE]}），本文保留该标注。


反过来说，这也解释了 6.2 里实验 B 的分解为什么成立：E − B2 ≈ 105 ms 是桥本体（可被特化消灭），B2 − E′ ≈ 16 ms 是影子栈内存纪律（手工替换消灭不了，但源头修复可以）。patch 把两者一起处理掉。


\subsection{6.5 两条路的取舍}

\begin{center}\small\begin{tabular}{lrrrrr}
\toprule
手段 & 实测 P50 (ms) & vs E & 零上游依赖 & 实现工作量 & 风险 \\
\midrule
E：perry 原样（未 pass 基线） & 123.098 & 1× & 是 & 0 & — \\
A5：wasm-opt 强内联 & 93.497 & 0.76× & 是 & \textasciitilde{}1 h & 低（体积 74 KB → 1.35 MB aot） \\
wamrc \texttt{--enable-segue} & 99.84 & 0.81× & 是 & \textasciitilde{}0.5 h & 仅 linux x86-64（GS 基址每线程寄存器） \\
通用桥内联 pass & \textbf{16.958} & \textbf{0.138×（快 7.3×）} & 是 & \textbf{\textasciitilde{}10 h} & 中：依赖 perry 产物指令形态，perry 升级 codegen 即失配 \\
本 pass + wasm-opt(A5) & 16.034 & 0.130× & 是 & +0 & 低 \\
本 pass + segue & 17.315 & 0.141× & 是 & +0 & 低 \\
B2：等价手工特化 & 17.185 & 0.141× & 是（但不可复用） & \textasciitilde{}4 h & 上界估计器，非修复 \\
V3：纯 f64 全去盒 & 3.325 & 0.027× & 是（但不可复用） & \textasciitilde{}8 h & 上界估计器，非修复 \\
\textbf{上游 patch（codegen 发射点特化）} & \textbf{3.891} & \textbf{0.0319×（快 31.4×）} & 否 & 未单独计量 & 见 7.2、6.3 遗留风险 \\
E′：干净 i64 wasm（引擎同速锚） & 1.216 & 0.010× & 是 & — & — \\
\bottomrule
\end{tabular}\end{center}


叠加结论 [S2:857-859]：wasm-opt 叠在 pass 之上只剩 −5\%（16.96 → 16.03），\texttt{--converge} 等强组合没有额外收益；\textbf{segue 在桥被消灭后不再有用}（16.96 → 17.32，方向反转），它的收益本来就来自优化 AOT 里那条桥分派路径，而 pass 已把该路径删掉。三者叠加无意义。


\textbf{核心回答}：不改 perry 上游，能把 122 ms 拉到 ≤20 ms 量级（16.0–17.0 ms，即 B2 上界水平），最小手段是单个 wat→wat 后处理 pass，接在既有 E 路链路的 \texttt{patch\_merged} 之后、\texttt{wasm-as} 之前，构建链只多一行命令；代价是约 10 小时一次性投入加随 perry 版本回归的风险 [S2:861-864]。改了上游，则直接到 3.891 ms，此时后处理 pass 可以整体退役：codegen 发射点特化是"源头修"，产物更紧，且不依赖 wat 后处理基础设施 [S7:103-108]。


\bigskip\hrule\bigskip

\section{7. 工程化可复用性讨论}

\subsection{7.1 方案层为什么正确}

三种修复手段（wasm-opt 后处理、wat 通用 pass、上游发射点特化）里，只有最后一种是可长期持有的，理由不在"快多少"，而在方案层性质。


\textbf{编译期特化，而非运行时优化。} 桥调用问题的本质是"信息在编译期就存在，却在运行期才被恢复"。\texttt{js\_add(number, number)} 的语义在发射点即可判定，perry 的原生后端已经这样做了（\texttt{typed\_abi.rs}、\texttt{i32\_fast\_path.rs}、\texttt{Type::Int32}）。把这个判定挪回发射点，等于把运行时的一次动态分派换成编译期的一次分支，问题在产生它的层次被消掉。wasm-opt 的失败恰好反证了这一点：它能内联整条桥，却无法折叠分派，因为分派表的索引来自运行时的内存 load，binaryen 没有内存常量传播 [S2:730-738]。


\textbf{等价性可证。} number 的表示是裸 f64 位型，\texttt{js\_add} 在 number 上的行为就是 f64 加法，内联后的 \texttt{i64.reinterpret\_f64(f64.add(...))} 与之逐位相同；盒布尔的 \texttt{is\_truthy} 就是 \texttt{!= TAG\_FALSE}。影子栈方面，原路径的 sp 净增减为 0，特化路径不碰 sp，净效果一致。这不是"看起来一样"，而是可以逐位论证的等价性 [S7:41-64]。不可证的地方一律回退：字符串 \texttt{+}、number 条件、\texttt{Mod}/\texttt{Pow}、\texttt{Eq}/\texttt{Ne}、闭包捕获、跨 module 函数返回值特化，全部保留原桥，行为与基线逐字节一致 [S7:89-92, S7:144-152]。


\textbf{向"上游的既有方向"收敛。} 上游已有 \texttt{--opt-report} 对特化拒绝原因的完整分类，说明类型特化是既定路线 [S2:437-438]；HIR 侧的类型推断（\texttt{infer\_expr\_type} / \texttt{infer\_binary\_type}）现成可调用，无需新写推断 [S2:792-795]。因此这一 patch 属于"补齐 wasm 后端与原生后端之间的落差"，而不是引入一套新的设计。


\subsection{7.2 缺口}

\textbf{自写类型推断应改接 HIR。} patch 里的 \texttt{type\_facts.rs} 是一个保守的、跨函数能力有限的轻量推断器。它之所以必要，是因为 wasm 后端至今没有把 HIR 的类型环境接进来。正确形态是阶段 0 的"装配 HIR 类型环境"，让 \texttt{WasmModuleEmitter} 持有 \texttt{HirTypeEnv}，发射点查询它而不是查询自建事实表 [S4:75-90]。自建推断的价值在于快速验证收益（这一轮实测证明了 31.4×），但它不该作为长期形态存在。


\textbf{测试面不足。} 现有正确性证据是三块：基准输出的逐字节校验、\texttt{demo.sh} 6/6（含一个负向用例）、三个泛化探针。这在"证明 patch 没把基准改坏"上是充分的，在"证明 patch 对所有 TS 写法都安全"上不是。回退路径的存在降低了风险（不可证即回退，行为等于基线），但回退触发率、回退路径本身的回归面都没有测量 [S7:144-152]。


\textbf{版本跟随成本。} 三条线都要随 perry 上游演进：


\begin{enumerate}
\item wat 后处理 pass 依赖 perry 产物的\textbf{指令形态}（\texttt{(drop (call \$mem\_call (f64.const <nameId>) …))} 加影子栈槽位约定），perry 一改 codegen 或桥发射形态就可能失配，这是后处理相对上游改造的\textbf{永久劣势}，必须随 perry 版本回归 [S2:830-832]。
\item 上游 patch 以 diff 形式存在，需要随上游 rebase；\texttt{type\_facts.rs} 依赖 \texttt{walker::walk\_expr\_children} 的穷尽匹配，上游新增 Expr 变体时会在编译期报错（这算是好事，属编译期而非运行期失效）。
\item typed ABI 化的长期路线（第 7.4 节引用 \texttt{docs/typed-abi-migration.md}）一旦落地，本文 patch 的发射点分支需要重新对齐到 typed 值表示。
\end{enumerate}

\textbf{上游 patch 自身的遗留风险}（原文照录）[S7:144-152, S2:1074-1079]：字符串 \texttt{+}、number 条件、\texttt{Mod}/\texttt{Pow}、\texttt{Eq}/\texttt{Ne} 仍走桥（正确性所需或保守回退）；类方法体/闭包体内无注解局部、跨 module 导入函数返回类型未纳入数据流，因此保守回退（无收益但无误判）；类型注解沿用 perry 上游语义（与原生 typed ABI 同样信任声明类型）。


\subsection{7.3 长期路线：typed ABI 化}

当前 patch 跨过了 B2 天花板，但与干净 i64 wasm（E′ 1.216 ms）之间还有一截，来源是影子栈内存纪律。要去掉它，需要把 wasm 后端的值表示与调用纪律整体改为 typed。这是已规划的"路径 4"，其六阶段方案见 \texttt{docs/typed-abi-migration.md} [S4]：


\begin{center}\small\begin{tabular}{lrrr}
\toprule
阶段 & 内容 & 验收 P50 & 人日 \\
\midrule
0 & 装配 HIR 类型环境（\texttt{HirTypeEnv::from\_module} 挂到 \texttt{WasmModuleEmitter}） & \textasciitilde{}122（无回归） & 0.5–1 \\
1 & 发射点特化（\texttt{+} / 条件内联）= 路径 3 完整 & \textasciitilde{}17.2（锚定 B2） & 2–4 \\
2 & 字面量/局部 typed（\texttt{INT32\_TAG} 0x7FFE + 局部 widening） & \textasciitilde{}12–17（推断） & 2–4 \\
3 & 签名 typed + trampoline + 拒绝制 & \textasciitilde{}12–17（推断） & 4–7 \\
4 & 去影子栈（typed 函数溢出改 local） & \textasciitilde{}3.3–5.0（V3 锚） & 4–8 \\
5 & \texttt{rt.*} 桥 typed 重载双轨 & \textasciitilde{}3.3 & 3–6 \\
\bottomrule
\end{tabular}\end{center}


合计约 16–30 人日（一人约 3–6 周）；两个可发布里程碑：M1（阶段 1 后 = B2，可发布）、M2（阶段 4 后 = V3，可发布）[S4:296-318]。方案里有一条非显见的结论：\textbf{阶段 2、3 对本基准的直接收益近零}，因为 fib 已是 TS 注解函数、递归调用已是直接 wasm \texttt{Call}、算术已在阶段 1 内联；它们真正的价值是作为阶段 4 的前置基础设施，14 ms 的跃迁发生在阶段 4 [S4:30-38]。


需要说明的是，本文实验（6.3）只实现了阶段 0 与阶段 1 的合并形态，且实现方式与规划文档不同（自建 \texttt{type\_facts} 而非装配 HIR 类型环境）。规划文档中阶段 2–5 的收益均为规划值，其中带推断标注的两行尚未实测。


\subsection{7.4 待补实证}

在本文写作时，一项针对"工程化可复用性"的独立实证工作正在并行进行，预计补入 \texttt{docs/performance.md} 的新增章节，内容包括：跨 perry 版本的 patch 可应用性探测、面向类/闭包/数组写法的泛化探针矩阵、以及 HIR 类型环境在上游 main 分支上的可用性核对。本文写作时尚无该章节落库，因此以上三节（7.1–7.3）的结论只基于 \texttt{tools/attribution/patch\_notes.md} 与 \texttt{docs/typed-abi-migration.md} 的如实声明，未包含该并行工作的实测结果。该章节落地后，第 7 章应据其更新，特别是"测试面不足"与"版本跟随成本"两条。


\bigskip\hrule\bigskip

\section{8. 局限与效度威胁}

\subsection{8.1 基准只覆盖一种负载}

全文的性能结论都建立在单个基准上：\texttt{fib(29)} 加 10⁶ 次求和循环。这是\textbf{调用密集的最坏情形}，它把跨边界调用成本放到最大，也正因如此，文中所有"倍数"都应当被理解为这一形态下的倍数，而不是 wasm 路线的普遍性能。


归因分析自己给出了限定词 [S2:440-444]：优化前的 79×/253× codegen 因子等于"可内联而未内联的算术/条件过桥（缺陷，可修复）"加"桥实现低效（nameId 线性扫描，实测值 34×）"加"NaN-box 指令形态"三者的乘积；其中"必要跨边界"（字符串、console 等）的成本不含在这三个因子里，基准未覆盖这类负载，对字符串密集负载应理解为"必要的跨边界成本加同样的实现低效"。


这一缺口在 6.2 得到部分弥补：\texttt{probe\_str.ts}（字符串密集）与 \texttt{probe\_mixed.ts}（混合类型）进入了泛化验证，收益分别是 2.0× 与 3.6×。但它们只用于验证 pass 的覆盖率与误判，没有被纳入第 4、5 章的乘性分解矩阵。\textbf{字符串密集负载在 A/E 两路中的占比仍未测量}，文档已将其列为限制 [S2:480-481]。


\subsection{8.2 测量环境与噪声}

机器是共享宿主上的 PVE 虚拟机，背景负载未做隔离，同机重测的 A′ 干净基线本身散布 50.8–58.8 ms，噪声约 ±15\% [S2:475-479]。具体影响：


\begin{itemize}
\item E 批间漂移：正式批 146.829 ms，复测批 148.709 / 138.736 ms（±7\%）[S2:485-487]；
\item D 四批复测 P50：6.651 / 6.558 / 6.428 / 6.286 ms（散布 ±6\%）[S2:545]；
\item B 扣除的 node 启动基线（优化前 23 ms、优化后 24 ms）在 1.24–1.28 s 里占约 2\%，不确定度与之同量级 [S2:475-476]。
\end{itemize}

\texttt{perf} 未安装，改用 callgrind（指令数精确，周期数由 P50 × 实测频率约 3.0 GHz 推算，±15\% 噪声）[S2:210-211]。另一处工具限制：valgrind 下 WAMR 宿主因 \texttt{touch\_pages} 栈增长受限而 SIGSEGV（8M/128M/512M 主栈同样），A/A′ 路的动态指令数无法用 callgrind 直接测，也没能给出"wasm 路每语义步指令数 vs 原生"的直接比值；该比值由 35× 与各自 IPC 估算，属推断，估计 wasm 路指令量约为原生的 10–20 倍 [S2:224-227]。


\subsection{8.3 口径不一致是刻意保留的}

六条路的计时方式不统一，这是产物形态决定的，不是疏忽：D 含进程启动（口径对 D 偏保守，若同口径扣启动则 D ≈ 5.3–5.7 ms，D/C ≈ 3.8–4.1×）[S2:534-538]；E 每轮独立进程（对 E 略偏不利，每轮多一次进程内预热但无进程启动成本，EXECUTE 计时段不含）[S2:482-485]；A 的冷启动 real 含 \texttt{bench\_time} 固定的一次预热，与 B/C 的"一次执行"口径不同，已单列 [S2:477-478]。


由此派生一条重要限定：\textbf{"23×"是 E 与 D 的进程级口径之比，不是同等口径之比}。按 D 的纯执行口径换算，E/D 是 54–88× [S2:610-617]。本文在两处都给出，读者应优先使用后者做算法层面的比较。


\subsection{8.4 demo 只覆盖语言的一个子集}

运行时的 13 个实现覆盖字符串、number/bool、console 与动态分派。\textbf{对象、数组、闭包、类一概没有}：那需要把 perry 的 handle store 在运行时模块里重做一遍，包括原型链、属性查找与 GC 语义 [S1:233]。198 个桩把边界标得很清楚，但也意味着本文的"到处运行"结论只对使用原始值子集的程序成立。


补齐这一子集的量级不是"再加几十个函数"：把 perry 的运行时语义手写重做一遍（无论 C 还是 Rust），等于重写 \texttt{perry-runtime}，还得跟着上游 ABI 走 [S1:237]。路线三把"扩展方式"从"写宿主"变成"在 \texttt{runtime-wasm/src/lib.rs} 里加实现、重跑 \texttt{gen-rt-symbols.mjs}" [S1:235]，但没有改变补齐完整 JS 语义所需的量级。适合这套方案的是\textbf{宿主可控、语言子集可裁剪}的场景（嵌入式规则脚本、计算密集的插件、既不想源码外流又不想放弃 TS 写法的内部交付）；把用满 npm 生态的应用搬过来，首先要解决的问题不是源码保护，而是运行时要补多少 [S1:249]。


\subsection{8.5 上游处于快速迭代期}

perry 的行号与内部结构都在变动。本文引用的上游行号分别锚定在 commit \texttt{87ecb02b}（patch 目标）[S7:5] 与另一份核对记录（HEAD \texttt{7ac11b09}）[S4:5-7]，两份 commit 并不相同，因此文中 file:line 只能作为"该版本下的位置"，不能作为长期稳定的坐标。typed ABI 规划文档也明确说明其全部 file:line 已对 \texttt{/tmp/perry-src} HEAD \texttt{7ac11b09} 逐行复核 [S3:455-456]。


同理，patch 的实验环境是 vendored 源码区（\texttt{/tmp/typerry-src/perry}，未 commit），基线由 tarball 解压后 \texttt{git init} 建立（commit \texttt{b4fef3e}），diff 由 \texttt{git diff} 生成 [S7:8, S3:512-514]。这意味着复现需要在同样的 commit 上做。


\subsection{8.6 推断性数字清单}

以下结论在原文中已标注推断，本文保留同等不确定性，汇总见附录 E。阅读本文数字时应注意的划分：


\begin{itemize}
\item \textbf{实测}：六路 P50 表、乘性因子与闭合校验、基线审计三条证据、隔离实验 \texttt{nohost} 系列、patch 前后 P50 与反汇编计数、pass 的 28 次逐字节一致与覆盖率。
\item \textbf{推断}：rt 分派/装箱占每层 1.2 µs 的大头（未逐项插桩）；B 路 364× 内 codegen 形态与 JS 宿主层的相对占比未拆分；E 路 49×→108× 的分母定义说明；D 路 fib/循环拆分（约 4 ms / 约 1.5 ms）由指令量比例推算；D 编译 1.9 s 的耗时构成未拆分；B2 的 reinterpret 对在 LLVM 下为空操作；V3 与 E′ 的 2.1 ms 差为 LLVM 对 f64 与 i64 的内联差异；pass 实现的约 10 小时为投入量级估计而非受控工时测量；patch 快于 B2 的 4.4× 归因于"整条帧建立与内存槽往返不再发射"。
\item \textbf{条件性 / 未复现}：坑 4 的 \texttt{initializing thread failed!}；\texttt{--enable-llvm-pgo} 未验证。
\item \textbf{未测}：WAMR LLVM JIT（需 \texttt{build\_llvm.sh} 自编全量 LLVM，收益与 AOT 同源，暂无必要）[S2:479-480]；E 路冷启动 [S2:485]；字符串密集负载在 A/E 两路中的占比 [S2:480-481]。
\end{itemize}

\bigskip\hrule\bigskip

\section{9. 结论}

回到引言的两个问题。


\textbf{RQ1（源码保护）大体成立，但带折扣。} 产物里没有 name 段，函数名与局部变量名都不泄漏；泄漏的是字符串字面量与导入名，后者还标出了程序会碰哪些运行时能力。它把门槛从"打开源码"抬到"反编译一遍再读"，真要用它保住商业逻辑，该藏的还得另外藏 [S1:241-243]。


\textbf{RQ2（到处运行还剩多少）已经有了明确答案：取决于运行时放在哪一侧。} 把运行时语义留给宿主，等价于要求每个平台重写一套运行时，成本随程序用到的语言特性线性增长（路线一的 C 宿主实测：13 个实现只够纯原始值程序，换成数组立刻报错）。把 perry 的 Rust 运行时按 wasm 目标编译成独立模块、由多模块机制链接，宿主就只剩 WASI 的一个调用（\texttt{fd\_write}），一次分发这一半成立。代价是运行时的 OS 强耦合模块需要裁剪，且完整 JS 语义（对象、闭包、GC）补齐的量级不变。


性能方面，比数字更值得留下的一条规则是：总倍数必须读作乘性因子的乘积。解释器下 1757× = 引擎因子 35× × codegen 因子 49×（闭合误差 3.1\%）；换 AOT 引擎后 104× = 0.97× × 108×（闭合误差 <1\%）。引擎因子在 AOT 下归零，干净 wasm 与手写 C 同速。\textbf{根因不是引擎，而是 perry wasm codegen 的类型擦除}：\texttt{+} 与条件判定被强制过桥，桥函数体占 AOT 路耗时的 95.8\%，每桥 25.4 ns。三处独立证据支持这一结论：隔离实验（与 perry 相同调用图加平凡被调方只需 1.371 ms）、V8 的 TurboFan 稳态（3.400 ms）、以及 QuickJS 旁证（纯解释器 85.532 ms 比 AOT 执行 perry 装箱字节码还快 0.58×）。


修复方面，两条路都走通了并各自量化。零上游依赖的 wat 后处理 pass 把 bench 从 123.1 ms 压到 16.9 ms（快 7.3×），4 程序 × 7 变体共 28 次逐字节一致、零误判，代价是约 800 行 JS 与随版本回归的维护负担；它的天花板是方法固有的：类型知识在模块里不可恢复时只能保守拒绝（\texttt{probe\_nested} 保守模式 43\% 对 \texttt{--closed-world} 86\%）。上游 patch 改 \texttt{perry-codegen-wasm} 的两个发射点并加一份保守类型事实，把 122.112 ms 降到 \textbf{3.891 ms（快 31.4×）}，跨过手工特化上界 17.185 ms，逼近全去盒的 3.325 ms，且 \texttt{demo.sh} 6/6、探针逐字节一致。它比手工特化还快 4.4× 的原因，是整条帧建立与内存槽往返都不再发射，而不只是替换了桥调用本身（该归因标注为推断）。


方法学上改动最小的一步是\textbf{审计自己的基线}。质疑"1.471 ms 物理上不可能"时，正确的回应不是复测一遍，而是查清 1,664,079 次逻辑调用里只有 91,759 次真实 call（gdb 断点与 callgrind 双证）。这条修正没有改变任何数字，只改变了 1757× 的解释方式：基准里的"调用次数"未必是硬件看到的调用次数。


\bigskip\hrule\bigskip

\section{参考资料}

本文不引用外部文献。全部引用为仓库内一手实测记录与上游公开仓库。


\textbf{仓库内材料}


\begin{itemize}
\item {[S1]} \texttt{docs/perry-wasm-wamr.md} —《把 TypeScript 编译成 wasm，塞进 WAMR 里跑起来》架构与 ABI 逆向记录（258 行）。
\item {[S2]} \texttt{docs/performance.md} —《perry wasm demo 性能基准》六路基准、归因分析、基线审计、修复路径与天花板（1083 行）。
\item {[S3]} \texttt{docs/process.md} —《从 C 桥接到 wasm 运行时：实施过程记录》背景、决策、实施步骤、踩坑、验证与后续增补（544 行）。
\item {[S4]} \texttt{docs/typed-abi-migration.md} —《路径 4：perry-codegen-wasm typed ABI 化 + 去影子栈 — 实施规划》（397 行）。
\item {[S5]} \texttt{README.md} — 项目总览、目录结构、实测输出、互操作约定与约束。
\item {[S6]} \texttt{tools/attribution/README.md} — 归因分析复现指南、实验文件清单与结果表。
\item {[S7]} \texttt{tools/attribution/patch\_notes.md} — 上游 patch 的设计说明、逐处语义与等价性论证、实测结果与遗留风险。
\item {[S8]} \texttt{src/bench.ts} — 基准程序（20 行）。
\end{itemize}

\textbf{上游仓库与产物}


\begin{itemize}
\item {[S9]} \href{https://github.com/PerryTS/perry}{PerryTS/perry} — perry 编译器（Rust，SWC + LLVM）。patch 目标 commit \texttt{87ecb02b}；行号核对 commit \texttt{7ac11b09}。
\item {[S10]} \href{https://github.com/fn-a/typerry}{fn-a/typerry} — perry wasm 后端的 npm 包 \texttt{@typerry/node}（submodule \texttt{d13b5769}）。
\item {[S11]} \href{https://github.com/bytecodealliance/wasm-micro-runtime}{bytecodealliance/wasm-micro-runtime} — WAMR 2.4.3，多模块链接、AOT（\texttt{wamrc}）与 \texttt{fd\_write} 支持。
\item {[S12]} perry v0.5.1520 预编译 release：\texttt{https://github.com/PerryTS/perry/releases/download/v0.5.1520/perry-linux-x86\_64.tar.gz}（sha256 \texttt{3423d9fea9bce9b2011fa53b5788a5ca115c947352b5c67278147de30fd2f952}）[S2:94-97]。
\item {[S13]} Bellard QuickJS（本地构建 \texttt{/tmp/quickjs-bellard}，\texttt{make qjs}）[S2:644-645]。
\item {[S14]} binaryen \texttt{wasm-opt} 123、\texttt{wasm-merge}、\texttt{wasm-dis}；WABT（\texttt{wasm2wat}、\texttt{wat2wasm}、\texttt{wasm-as}）[S2:106-115, S2:716]。
\end{itemize}

\bigskip\hrule\bigskip

\section{附录 A 完整数字表}

\subsection{A.1 六路稳态 P50（优化采纳后主批）[S2:136-146]}

\begin{center}\small\begin{tabular}{lrrrr}
\toprule
目标 & P50 & 最小 & 最大 & ÷C \\
\midrule
E. WAMR AOT（合并单模块） & 146.829 ms & 143.097 & 158.074 & 104.4× \\
E′. WAMR AOT（干净 wasm） & 1.362 ms & 1.317 & 1.409 & 0.97× \\
A. WAMR 解释器 & 2470.678 ms & 2378.375 & 2527.101 & 1757.2× \\
B. Node V8（扣启动 24 ms） & 1280.000 ms & 1246.000 & 1311.000 & 910.4× \\
C. 原生 gcc -O2 & 1.406 ms & 1.357 & 3.092 & 1× \\
D. perry 原生（进程级） & 6.286 ms & 5.946 & 7.306 & 4.5× \\
F. QuickJS（进程级） & 85.532 ms & — & — & 61× \\
\bottomrule
\end{tabular}\end{center}


\subsection{A.2 优化前对照 [S2:150-158]}

\begin{center}\small\begin{tabular}{lrrr}
\toprule
目标 & P50 & 最小 & 最大 \\
\midrule
A. WAMR 解释器（优化前） & 4009.746 ms & 3866.041 & 4499.434 \\
B. Node V8（扣启动 23 ms） & 1239.000 ms & 1204.000 & 1387.000 \\
C. 原生 gcc -O2 & 1.471 ms & 1.322 & 1.692 \\
\bottomrule
\end{tabular}\end{center}


倍数：A = 2725.9×，B = 842.3×。


\subsection{A.3 乘性分解与闭合校验}

\begin{center}\small\begin{tabular}{lrrrrr}
\toprule
路 & 引擎因子 & codegen 因子 & 乘积 & 实测 & 误差 \\
\midrule
A（解释器，优化后） & \textasciitilde{}35× & \textasciitilde{}49× & 1703 & 1757 & 3.1\% \\
A（解释器，优化前） & \textasciitilde{}35× & \textasciitilde{}79× & 2730 & 2726 & 0.2\% \\
B（V8，审计后） & \textasciitilde{}2.5× & \textasciitilde{}364× & 910 & 910.4 & <0.1\% \\
B（V8，审计前） & \textasciitilde{}3.3× & \textasciitilde{}253× & 835 & 842 & — \\
E（AOT） & \textasciitilde{}0.97× & \textasciitilde{}108× & 105 & 104.4 & <1\% \\
\bottomrule
\end{tabular}\end{center}


\subsection{A.4 E 路 AOT 闭合（两套等价分解）[S2:593-605, S3:392]}

\begin{center}\small\begin{tabular}{lrr}
\toprule
成分 & 实测 & 来源 \\
\midrule
fib 纯机器码基线 & 1.309 ms & \texttt{fib\_only.aot} \\
loop 纯机器码基线（LLVM 常量折叠） & \textasciitilde{}0.002 ms & \texttt{loop\_only.aot} \\
桥机制 + 装箱往返税（每层 2 次调用，最小 i64↔f64 往返） & 5.913 ms（税 4.604） & \texttt{nohost\_box\_app.wat} \\
rt 侧 \texttt{mem\_call} 函数体执行 & 135.5 ms & 余项 \\
\textbf{合计} & \textbf{141.4 ms} & 闭合 ✓（误差 <0.1\%） \\
\bottomrule
\end{tabular}\end{center}


另一写法：1.309 + 4.604 + 135.5 ≈ 141.4 ms，对实测 141.382 ms [S3:392]。每桥成本：rt 函数体 25.4 ns（÷ 5,328,158 桥）、装箱税 1.4 ns（÷ 3,328,158 桥，仅 fib）。


\subsection{A.5 基线审计三证据 [S2:194-216]}

\begin{center}\small\begin{tabular}{lrrrr}
\toprule
编译 & 单次指令数 & 静态 call 点 & P50 & IPC \\
\midrule
gcc -O1 & 28.13 M & 2 & 3.604 ms & \textasciitilde{}2.6 \\
gcc -O2（基线） & 16.71 M & 4（2 克隆体） & 1.471 ms & \textasciitilde{}3.8 \\
\bottomrule
\end{tabular}\end{center}


\begin{center}\small\begin{tabular}{lr}
\toprule
编译器/选项 & P50 \\
\midrule
gcc -O1 & 3.604 ms \\
gcc -O2 & 1.471 ms \\
gcc -O3 & 1.240 ms \\
gcc -Ofast & 1.296 ms \\
gcc -Os & 2.337 ms \\
clang -O2 & 2.170 ms \\
\bottomrule
\end{tabular}\end{center}


拆分计时（gcc -O2）：fib 0.84 ms + 循环 0.30 ms ≈ 1.13 ms（整体 1.47 ms）；同型 f64 变体 2.52 ms。


\subsection{A.6 修复实验对照（2026-09-21 同口径）}

\begin{center}\small\begin{tabular}{lrr}
\toprule
变体 & P50 (ms) & 相对 E \\
\midrule
E（perry 原样，patch 批） & 122.112 & 1× \\
A1 / A2 / A3 / A4 / A5（wasm-opt 变体） & 135.4 / 99.6 / 101.4 / 101.3 / 93.497 & ≈1× / 0.78–0.81× / 0.82× / 0.82× / 0.766× \\
B1 / B2 / V3 & 71.612 / 17.185 / 3.325 & 0.586× / 0.141× / 0.027× \\
E′（同批 clean i64） & 1.216 & 0.010× \\
pass 批 E 基线 & 123.098 & 1× \\
pass（bench） & 16.958 & 0.138× \\
\textbf{上游 patch} & \textbf{3.891} & \textbf{0.0319×} \\
\bottomrule
\end{tabular}\end{center}


\subsection{A.7 泛化验证 P50 全表（ms）[S2:931-936]}

\begin{center}\small\begin{tabular}{lrrrrrrr}
\toprule
程序 & base & pass & pass(cw) & pass+wasmopt & pass+wasmopt(强) & pass+segue & pass+wasmopt+segue \\
\midrule
bench & 123.098 & 16.958 & 16.822 & 16.034 & 16.203 & 17.315 & 17.661 \\
probe\_str & 30.297 & 14.942 & 14.340 & 10.213 & 14.248 & 13.429 & 11.908 \\
probe\_mixed & 21.261 & 5.918 & 5.861 & 4.421 & 5.629 & 5.058 & 4.690 \\
probe\_nested & 0.521 & 0.280 & 0.044 & 0.236 & 0.282 & 0.233 & 0.226 \\
\bottomrule
\end{tabular}\end{center}


\subsection{A.8 覆盖率与误判}

\begin{center}\small\begin{tabular}{lrrr}
\toprule
程序 & 改写前→后桥调用点 & 保守 & closed-world \\
\midrule
bench & 10 → 6 & 40\% & 40\% \\
probe\_str & 11 → 6 & 45\% & 45\% \\
probe\_mixed & 7 → 4 & 43\% & 43\% \\
probe\_nested & 7 → 4 & 43\% & 86\% \\
\bottomrule
\end{tabular}\end{center}


正确性：28/28 逐字节一致；误判清单为空。


\subsection{A.9 patch 前后反汇编计数 [S7:131-137]}

\begin{center}\small\begin{tabular}{lrr}
\toprule
指标 & 前 & 后 \\
\midrule
文件大小 & 9780 B & 9561 B \\
\texttt{call \$mem\_call} & 8 & 6 \\
\texttt{call \$mem\_call\_i32} & 2 & 0 \\
\texttt{f64.add} & 0 & 3 \\
\texttt{i64.ne} & 0 & 2 \\
\bottomrule
\end{tabular}\end{center}


\subsection{A.10 产物尺寸}

\begin{center}\small\begin{tabular}{lr}
\toprule
文件 & 字节 \\
\midrule
\texttt{build/app.wasm} & 10650 \\
\texttt{build/app\_link.wasm} & 10658 \\
\texttt{build/rt.wasm}（优化前 / 采纳 nameId 缓存后） & 16798 / 16928 \\
\texttt{build/perry\_link}（宿主 runner） & 531336 \\
\texttt{runtime-wasm/src/lib.rs} & 620 行 \\
\texttt{build/bench\_perry\_native}（perry 原生产物） & 16.3 MB（17,099,736 B） \\
\texttt{build/bench\_merged.aot}（E 路 AOT 产物） & 74 KB \\
\bottomrule
\end{tabular}\end{center}


\bigskip\hrule\bigskip

\section{附录 B 复现命令清单}

以下命令均在各来源文档中给出，按路分列。所有命令在项目根目录执行。


\begin{verbatim}
# A. WAMR 解释器（完整链路构建后）[S2:75-82]
node tools/build-wasm.mjs src/bench.ts --bare build/bench.wasm
node tools/gen-rt-symbols.mjs build/bench.wasm --rust runtime-wasm/src/lib.rs build/rt_symbols.rs
cargo build --release --target wasm32-unknown-unknown --manifest-path runtime-wasm/Cargo.toml
cp runtime-wasm/target/wasm32-unknown-unknown/release/perry_rt_wasm.wasm build/rt_bench.wasm
node tools/patch-app-memory.mjs build/bench.wasm build/bench_link.wasm
gcc -O2 -I .deps/wamr/core/iwasm/include -o build/bench_time ...
./build/bench_time build/bench_link.wasm build/rt_bench.wasm 1

# B. perry JS 宿主层（V8 参照）[S2:84-85]
node build/ref/run.mjs

# C. 原生基线 [S2:87-88]
gcc -O2 -o build/bench_native tools/bench_native.c && ./build/bench_native 1

# D. perry 原生（一键：下载+校验+编译+计时）[S2:90-100]
tools/attribution/bench_d.sh 11

# E. WAMR AOT（一键，含 E/E' 构建+校验+计时）[S2:102-117]
tools/attribution/aot_e.sh 12

# F. QuickJS 对照 [S6:91-93]
make -C /tmp/quickjs-bellard qjs
tools/attribution/bench_f.sh 11

# 基线审计 [S6:60-107]
node tools/attribution/run_node_steady.mjs
./build/xmod_time build/nohost_app.wasm build/nohost_triv.wasm 7 | grep RUN

# 修复实验 [S6:38-42]
tools/attribution/exp_wasmopt.sh 12       # 实验 A：wasm-opt 变体
tools/attribution/exp_specialize.sh 12    # 实验 B：B1/B2/V3 手工特化
tools/attribution/exp_postpass.sh all     # 通用桥内联 pass：4 程序 × 7 变体

# 反汇编证据 [S6:84-87]
wasm2wat --enable-all build/bench.wasm -o /tmp/bench.wat
grep -c 'call 209\|call 210' /tmp/bench.wat
\end{verbatim}

一键演示与全量基准：\texttt{./demo.sh}（6 步）、\texttt{./tools/bench.sh}（末尾打印对照表）[S5:29-31, S2:5-7]。


\bigskip\hrule\bigskip

\section{附录 C \texttt{rt.*} ABI 速查}

\textbf{值编码（NaN-boxing，i64 携带 f64 位模式，高 16 位为标签）}[S1:103-111, S3:60-66]


\begin{center}\small\begin{tabular}{lr}
\toprule
值 & 位模式 \\
\midrule
\texttt{undefined} / \texttt{null} / \texttt{false} / \texttt{true} & \texttt{0x7FFC…0001} \textasciitilde{} \texttt{0x7FFC…0004} \\
对象 / 数组 / 闭包（handle） & 高 16 位 \texttt{0x7FFD}，低 32 位 handle id \\
int32 快路径 & 高 16 位 \texttt{0x7FFE}（wasm 后端从不发射） \\
字符串 & 高 16 位 \texttt{0x7FFF}，低 32 位字符串表下标 \\
其他 & 普通 double 位模式 \\
盒布尔（特化用常量） & \texttt{TAG\_TRUE} 0x7FF8000000000004 / \texttt{TAG\_FALSE} 0x7FF8000000000003 [\texttt{S2:765-766}] \\
\bottomrule
\end{tabular}\end{center}


\textbf{动态调用协议}[S1:115]


\begin{verbatim}
mem_call(nameId: i32, argc: i32, base: i32) -> f64(占位 0.0)
  · 参数：以 u64 槽位写在业务线性内存 base + i*8 处（u64::from_le）
  · 返回值：写回 base
mem_call_i32(nameId: i32, argc: i32, base: i32) -> i32
  · 结果直接作为 i32 返回，不写内存
\end{verbatim}

\textbf{桥 nameId 语义}（来自 rt 固定桥表，与数据段字符串序一致，\texttt{id = 序 + 1}）[S2:900-904]


\begin{center}\small\begin{tabular}{lrrr}
\toprule
nameId & 桥 & nameId & 桥 \\
\midrule
4 & \texttt{console\_log} & 10 & \texttt{string\_len} \\
8 & \texttt{js\_add} & 12 & \texttt{is\_truthy} \\
9 & \texttt{string\_eq} & 13 & \texttt{js\_strict\_eq} \\
\bottomrule
\end{tabular}\end{center}


\textbf{字符串表契约}[S1:113, S3:72-80]：启动时按固定顺序逐个调用 \texttt{rt.string\_new(offset, len)} 注册字面量，宿主/运行时必须按同样顺序 append，下标即 id。运行时侧 \texttt{MAX\_STRINGS = 1024} 条、\texttt{ARENA\_SIZE = 64 KiB}，条目 \texttt{\{ptr: u32, len: u32, utf16: u32\}}，溢出直接 \texttt{fatal}。


\textbf{13 个已实现导出}[S5:77-79]：\texttt{string\_new}、\texttt{console\_log}/\texttt{console\_warn}/\texttt{console\_error}、\texttt{string\_concat}、\texttt{js\_add}、\texttt{string\_eq}、\texttt{js\_strict\_eq}、\texttt{is\_truthy}、\texttt{string\_len}、\texttt{jsvalue\_to\_string}、\texttt{mem\_call}、\texttt{mem\_call\_i32}。另有 \texttt{\_initialize}（reactor 入口，WAMR 对带 WASI 导入的模块要求导出）[S3:110-111]。


\textbf{perry 产物的桥发射边界}[S2:373-380]


\begin{center}\small\begin{tabular}{lr}
\toprule
操作 & 路径 \\
\midrule
\texttt{+}（js\_add） & 过桥 \\
\texttt{-} \texttt{*} \texttt{/} & 内联 f64.sub/mul/div \\
\texttt{<} \texttt{<=} \texttt{>} \texttt{>=} & 内联 f64.lt/le/gt/ge \\
if/while/for 条件 & 过桥（is\_truthy） \\
\texttt{===} / \texttt{==} & 过桥（js\_strict\_eq） \\
字符串操作 / console & 过桥（本来就必须） \\
\bottomrule
\end{tabular}\end{center}


\bigskip\hrule\bigskip

\section{附录 D 文件清单}

\textbf{运行时与宿主}


\begin{center}\small\begin{tabular}{lr}
\toprule
路径 & 作用 \\
\midrule
\texttt{runtime-wasm/Cargo.toml} & \texttt{cdylib}、\texttt{opt-level="s"}、\texttt{lto}、\texttt{panic="abort"}、\texttt{codegen-units=1}、\texttt{strip} \\
\texttt{runtime-wasm/src/lib.rs} & \texttt{\#![no\_std]} 运行时，620 行，13 个实现 + \texttt{include!} 桩表 \\
\texttt{runtime-wasm/.cargo/config.toml} & \texttt{--global-base=2097152} 等链接参数 \\
\texttt{host/perry\_link.c} & WAMR 多模块 runner（load/register/set\_wasi\_args/instantiate/execute） \\
\texttt{host/perry\_rt.c}、\texttt{host/perry\_abi.h} & 路线一 C 桥接探针（历史） \\
\bottomrule
\end{tabular}\end{center}


\textbf{工具}


\begin{center}\small\begin{tabular}{lr}
\toprule
路径 & 作用 \\
\midrule
\texttt{tools/build-wasm.mjs} & TS → wasm（走 \texttt{@typerry/node} 的 \texttt{wasmBare} / \texttt{wasmBoot}） \\
\texttt{tools/gen-rt-symbols.mjs} & 解析导入段，生成 \texttt{build/rt\_symbols.rs}（13 实现 + 198 桩） \\
\texttt{tools/patch-app-memory.mjs} & 业务模块 memory 段改 \texttt{import rt.memory} \\
\texttt{tools/inspect-wasm.mjs} & 产物结构检查 \\
\texttt{tools/bench.sh} & A/B/C 三路对照与输出一致性校验 \\
\texttt{tools/bench\_time.c} & A 路进程内计时包装 \\
\texttt{tools/bench\_native.c} & C 路基线（经 \texttt{volatile} 反折叠） \\
\texttt{demo.sh} & 6 步一键演示 \\
\bottomrule
\end{tabular}\end{center}


\textbf{归因与修复实验}（\texttt{tools/attribution/}）


\begin{center}\small\begin{tabular}{lr}
\toprule
类别 & 文件 \\
\midrule
干净对照与计时 & \texttt{clean\_bench.wat}、\texttt{clean\_time.c}、\texttt{run\_node.mjs}、\texttt{run\_node\_steady.mjs}、\texttt{bench\_split.c}、\texttt{bench\_f64.c}、\texttt{fib\_count.c}、\texttt{fib\_O2.asm} \\
跨模块与隔离 & \texttt{nohost\_app.wat}、\texttt{nohost\_triv.wat}、\texttt{nohost\_triv\_nomem.wat}、\texttt{nohost\_box\_app.wat}、\texttt{xmod\_time.c} \\
D/E/F 路 & \texttt{bench\_perry\_wrap.c}、\texttt{bench\_d.sh}、\texttt{aot\_time.c}、\texttt{aot\_e.sh}、\texttt{patch\_merged.mjs}、\texttt{bench\_exec\_wrap.c}、\texttt{bench\_quickjs.js}、\texttt{bench\_f.sh} \\
实验 A/B & \texttt{exp\_wasmopt.sh}、\texttt{spec\_patch.py}、\texttt{exp\_specialize.sh}、\texttt{specialized\_bench\_f64.wat} \\
后处理 pass & \texttt{bridge\_inline\_pass.mjs}、\texttt{exp\_postpass.sh}、\texttt{probe\_ref.mjs}、\texttt{probes/probe\_\{str,mixed,nested\}.ts} \\
上游 patch & \texttt{codegen\_specialize.patch}、\texttt{patch\_notes.md} \\
开关侦察 & \texttt{switch\_recon/}（\texttt{README.md} + \texttt{REPORT.md}） \\
\bottomrule
\end{tabular}\end{center}


\bigskip\hrule\bigskip

\section{附录 E \texttt{[INFERENCE]} 与不确定性标注清单}

\textbf{\texttt{[INFERENCE]}}


\begin{center}\small\begin{tabular}{lrr}
\toprule
\# & 内容 & 出处 \\
\midrule
1 & rt 侧分派 + NaN-box 装箱构成每层约 1.2 µs 的大头（未逐项插桩） & [S2:285] \\
2 & 49× → 108× 的因子变化是分母定义改变的结果，两因子不可直接相除比较 & [S2:313-314] \\
3 & wasm 路每语义步指令量约为原生的 10–20 倍（由 35× 与各自 IPC 估算） & [S2:226-227] \\
4 & D 路 fib/循环拆分（约 4 ms / 约 1.5 ms）由指令量比例推算，未单独插桩；D 编译 1.9 s 的耗时构成未拆分 & [S2:547-548] \\
5 & nameId 折叠成直接调桥的收益被夹在 A5 与 B1 之间（只省 dispatch） & [S2:743-745] \\
6 & B2 的 i64↔f64 reinterpret 对在 LLVM 是空操作、box 本身近零成本 & [S2:774] \\
7 & V3 与 E′ 的 2.1 ms 差是 LLVM 对 f64 与 i64 fib 的内联/优化差异 & [S2:778-779] \\
8 & pass 实现约 10 小时为本次 spike 的实际投入量级估计，非受控工时测量 & [S2:828-829, S2:989] \\
9 & perry codegen 不会把非 number 喂进 f64 域运算（依据是它自己无条件内联 F64Sub/Mul/Div） & [S2:984-985] \\
10 & 影子栈帧约定由 4 个程序的产物归纳，未见于上游文档 & [S2:986-988] \\
11 & \texttt{--closed-world} 的安全性依赖"宿主只调 \texttt{\_start}" & [S2:990] \\
12 & B2/V3 与 E′ 之间的差（影子栈内存纪律）沿用实验 B 的推断 & [S2:991] \\
13 & patch 快于 B2 4.4× 的主因是帧建立 + 内存槽往返不再发射 & [S7:119-121, S2:1038-1040] \\
14 & \texttt{Integer} 字面量与 \texttt{Int32} 声明在 wasm 后端一律发射为 f64 位型 & [S7:96-98] \\
15 & 影子栈净增减不变的论证基于 sp 配对的源码注释，未逐指令仿真验证 & [S7:99-101] \\
16 & typed 签名（I64→F64）在 AOT 下是同位宽 reinterpret 空操作；trampoline 经 wamrc 内联后近零开销 & [S4:34-35] \\
17 & 阶段 2/3 的验收 P50（\textasciitilde{}12–17 ms）为规划值 & [S4:23-24, S4:143-144] \\
18 & 坑 1 的 \texttt{wasm-abis=3} 依据用户记录，仓库现无该痕迹 & [S3:172-179] \\
19 & 为何 \texttt{mem\_call} 绕内存而非直接传 f64（NaN 位模式规范化风险）为作者判断，源码未解释 & [S1:115] \\
20 & 约 800 行 JS / 每个 pass 变体叠加强组合的收益判断 & [S2:828, S2:857-859] \\
\bottomrule
\end{tabular}\end{center}


\textbf{条件性 / 未复现}


\begin{center}\small\begin{tabular}{lr}
\toprule
内容 & 出处 \\
\midrule
坑 4 \texttt{initializing thread failed!}：本 WAMR 2.4.3 构建未开 WASI 线程支持，未复现 & [S3:196-203] \\
wamrc \texttt{--enable-llvm-pgo} 因缺 \texttt{WAMR\_BUILD\_STATIC\_PGO=1} 无法闭环，未验证 & [S2:968-969] \\
\bottomrule
\end{tabular}\end{center}


\textbf{未测}


\begin{center}\small\begin{tabular}{lr}
\toprule
内容 & 出处 \\
\midrule
WAMR LLVM JIT（需自编全量 LLVM，收益与 AOT 同源） & [S2:479-480] \\
E 路冷启动 & [S2:485] \\
字符串密集负载在 A/E 两路中的占比 & [S2:480-481] \\
B 路 364× 内 codegen 形态与 JS 宿主层的相对占比 & [S2:466-467, S6:133-134] \\
minhost（JS 最小桩跑 perry wasm）实验因桩语义复杂超时放弃 & [S6:133-134] \\
\bottomrule
\end{tabular}\end{center}


\end{document}
